\title{Επικουρικές Υπηρεσίες μη Σχετικές με τη Συχνότητα}

\section{Εισαγωγή}
\label{kef-02-theoria}
Οι αγορές ηλεκτρικής ενέργειας εξασφαλίζουν ότι η προγραμματισμένη παραγωγή καλύπτει την
αναμενόμενη ζήτηση (Κεφάλαιο 1, Ενότητα «\nameref{sec:stages}»). Για την ασφαλή λειτουργία του
ΣΜΗΕ ο ΔΣΜΗΕ χρειάζεται όμως και \emph{επικουρικές υπηρεσίες} (ancillary services), δηλαδή
υπηρεσίες που διατηρούν τη συχνότητα και τις τάσεις εντός ορίων, στηρίζουν το σύστημα κατά τις
διαταραχές και επιτρέπουν την επανατροφοδότησή του μετά από γενική διακοπή. Η Οδηγία (ΕΕ)
2019/944 για την εσωτερική αγορά ηλεκτρικής ενέργειας τις διακρίνει σε υπηρεσίες εξισορρόπησης,
δηλαδή τις επικουρικές υπηρεσίες σχετικές με τη συχνότητα, και σε \emph{επικουρικές υπηρεσίες
μη σχετικές με τη συχνότητα} (non-frequency ancillary services). Η διάκριση αντανακλά τη
διαφορετική φύση των δύο μεγεθών. Η συχνότητα της τάσης είναι κοινή σε ολόκληρη τη
συγχρονισμένη περιοχή και σχετίζεται άμεσα με το ισοζύγιο ενεργής ισχύος, οπότε κάθε μονάδα της
περιοχής συμβάλλει στη ρύθμισή της· οι σχετικές υπηρεσίες εξετάζονται στο
\hyperref[kef-03-theoria]{Κεφάλαιο 3}. Στον αντίποδα, το μέτρο της τάσης διαφέρει από ζυγό σε
ζυγό και καθορίζεται κυρίως από τις τοπικές συνθήκες φόρτισης του δικτύου. Οι επικουρικές
υπηρεσίες μη σχετικές με τη συχνότητα παρουσιάζονται στον Πίνακα~\ref{tab:nfas}.

\begin{table}[!htbp]
\centering
\begin{tabular}{lll}
\hline
\textbf{Υπηρεσία} & \textbf{Σκοπός} & \textbf{Χρονική Κλίμακα} \\
\hline
Ρύθμιση τάσης σε μόνιμη κατάσταση & τάσεις εντός ορίων & s -- h \\
Γρήγορη έγχυση άεργου ρεύματος & στήριξη τάσης σε σφάλματα & ms -- s \\
Αδράνεια για την τοπική ευστάθεια & συγκράτηση συχνότητας & ms -- s \\
Ρεύμα βραχυκύκλωσης & λειτουργία προστασίας & ms \\
Ικανότητα επανεκκίνησης & αποκατάσταση συστήματος & min -- h \\
Ικανότητα νησιδοποιημένης λειτουργίας & τροφοδότηση νησίδας & min -- h \\
\hline
\end{tabular}
\caption{Επικουρικές υπηρεσίες μη σχετικές με τη συχνότητα, σύμφωνα με την Οδηγία (ΕΕ)
2019/944.}
\label{tab:nfas}
\end{table}

Στο παραδοσιακό ΣΗΕ οι υπηρεσίες αυτές παρέχονταν εγγενώς από τις σύγχρονες γεννήτριες
των μεγάλων σταθμών παραγωγής ηλεκτρικής ενέργειας. Η αντικατάστασή τους από μονάδες ΑΠΕ που συνδέονται μέσω
μετατροπέων ηλεκτρονικών ισχύος περιορίζει τη διαθεσιμότητα των παραπάνω υπηρεσιών: ένας μετατροπέας μπορεί να παρέχει αντίστοιχες
υπηρεσίες μόνο εφόσον σχεδιαστεί και ρυθμιστεί κατάλληλα. Για τον λόγο αυτό ο Κανονισμός (ΕΕ) 2016/631 για τις απαιτήσεις σύνδεσης
των μονάδων ηλεκτροπαραγωγής (Requirements for Generators, RfG) ορίζει ικανότητες άεργης
ισχύος, ρύθμισης τάσης και στήριξης τάσης σε σφάλματα, ενώ οι ΔΣΜΗΕ αναγκάζονται να εγκαταστήσουν
επιπλέον εξοπλισμό για την αντιμετώπιση του παραπάνω προβλήματος.

Από τις υπηρεσίες του Πίνακα~\ref{tab:nfas}, το Κεφάλαιο εστιάζει στη ρύθμιση της τάσης σε
μόνιμη κατάσταση, δηλαδή στα μέσα ρύθμισης, στην οργάνωση της ρύθμισης από τον προγραμματισμό
έως τον πραγματικό χρόνο και στην ευστάθεια τάσης, καθώς και στη στήριξη τάσης
σε σφάλματα. Ακολουθούν η ικανότητα επανεκκίνησης και η ικανότητα νησιδοποιημένης
λειτουργίας. Η αδράνεια για την τοπική ευστάθεια, αν και περιλαμβάνεται στον Πίνακα~\ref{tab:nfas}, σχετίζεται
άμεσα με τη συχνότητα, αφού καθορίζει τον αρχικό ρυθμό μεταβολής της μετά από μια διαταραχή του
ισοζυγίου ενεργής ισχύος, και γι' αυτό εξετάζεται στο \hyperref[kef-03-theoria]{Κεφάλαιο 3}.
Στα Παραδείγματα που ακολουθούν, οι έννοιες του κεφαλαίου εφαρμόζονται με κώδικα που εκτελείται
απευθείας στην ηλεκτρονική έκδοση.

\section{Ρύθμιση τάσης σε μόνιμη κατάσταση}
\label{sec:v-control}

Σκοπός της ρύθμισης τάσης είναι να διατηρούνται οι τάσεις όλων των ζυγών εντός των ορίων
επιχειρησιακής ασφάλειας, τόσο στην κατάσταση N όσο και μετά από κάθε απρόβλεπτο συμβάν του
καταλόγου (Κεφάλαιο 1, Ενότητα «\nameref{sec:n1}»). Παράλληλα, η ρύθμιση επιδιώκει τον περιορισμό των ροών άεργης ισχύος, ώστε να
μειώνονται οι απώλειες και να αποδεσμεύεται μεταφορική ικανότητα για την ενεργή ισχύ, καθώς
και τη διατήρηση επαρκούς εφεδρείας άεργης ισχύος για την αντιμετώπιση απρόβλεπτων
συμβάντων.

Στην Ηπειρωτική Ευρώπη, ο SOGL ορίζει ως εύρος τάσης της κανονικής κατάστασης
0,90--1,05 p.u. για τάσεις από 300 έως 400 kV και 0,90--1,118 p.u. για τάσεις από 110 έως
300 kV. Ο ΑΔΜΗΕ υιοθετεί τα εύρη αυτά στον Κώδικα Διαχείρισης του ΕΣΜΗΕ ως επιτρεπόμενο εύρος
τάσης των ζυγών των 400 kV και των 150 kV αντίστοιχα και παρακολουθεί τις πιθανές αποκλίσεις καταγράφοντας τη διάρκεια και τον μέγεθός τους για κάθε επίπεδο τάσης. Οι αποκλίσεις αυτές
κατατάσσονται κατά μέγεθος και διάρκεια στις κλίμακες της μεθοδολογίας ταξινόμησης συμβάντων
του ENTSO-E (Incident Classification Scale, ICS), όπως συνοψίζει ο Πίνακας~\ref{tab:v-limits}:
όσο μεγαλύτερη και διαρκέστερη είναι μια απόκλιση, τόσο σοβαρότερο θεωρείται το συμβάν.

\begin{table}[!htbp]
\centering
\begin{tabular}{lllll}
\hline
\textbf{Ζώνη τάσης} & \textbf{150 kV (p.u.)} & \textbf{400 kV (p.u.)} & \textbf{Διάρκεια}
& \textbf{Κατάταξη} \\
\hline
Εντός ορίων & 0,90--1,118 & 0,90--1,05 & απεριόριστη & -- \\
\multirow{3}{*}{Μικρή απόκλιση} & \multirow{3}{*}{0,85--0,90 ή\newline 1,118--1,15}
& \multirow{3}{*}{0,85--0,90 ή\newline 1,05--1,10} & 5--15 min & κάτω από την κλίμακα \\
 & & & 15--60 min & κλίμακα 0 \\
 & & & $>$ 60 min & κλίμακα 1 \\
Μεγάλη απόκλιση & $<$ 0,85 ή $>$ 1,15 & $<$ 0,85 ή $>$ 1,10 & $>$ 30 s & κλίμακα 1 \\
\hline
\end{tabular}
\caption{Επιτρεπτά όρια τάσης των ζυγών του ΕΣΜΗΕ και κατάταξη των αποκλίσεων σύμφωνα με τον Κώδικα
Διαχείρισης του ΕΣΜΗΕ. Ένα συμβάν κατατάσσεται στην κλίμακα 2 όταν,
με τα κριτήρια της κλίμακας 1, έχει επιπλέον συνέπειες σε γειτονικό ΔΣΜΗΕ.}
\label{tab:v-limits}
\end{table}

\subsection{Τάση και άεργη ισχύς}
\label{sec:v-q}

Η σχέση της τάσης με την άεργη ισχύ προκύπτει από τη γραμμή δύο ζυγών του
Σχήματος~\ref{fig:vq}, με σύνθετη αντίσταση σειράς $R + jX$, μέσω της οποίας ρέει ισχύς
$P + jQ$ προς τον ζυγό 2. Αν ο φασιθέτης $\tilde{V}_2 = V_2$ ληφθεί ως αναφορά, το ρεύμα της
γραμμής είναι $\tilde{I} = (P - jQ)/V_2$ και η τάση του ζυγού 1 είναι

\begin{equation}
\tilde{V}_1 = V_2 + (R + jX)\,\tilde{I}
= V_2 + \frac{RP + XQ}{V_2} + j\,\frac{XP - RQ}{V_2}
\label{eq:v1phasor}
\end{equation}

Ο δεύτερος όρος βρίσκεται σε φάση με τον $\tilde{V}_2$ και καθορίζει κυρίως τη διαφορά των
μέτρων των τάσεων, ενώ ο τρίτος είναι κάθετος και καθορίζει τη διαφορά των γωνιών
$\delta = \theta_1 - \theta_2$. Όταν η γωνία $\delta$ είναι μικρή, ισχύει κατά προσέγγιση

\begin{equation}
\Delta V = V_1 - V_2 \approx \frac{RP + XQ}{V_2}, \qquad
\delta \approx \frac{XP - RQ}{V_2^2}
\label{eq:dv}
\end{equation}

Στα δίκτυα μεταφοράς η αντίδραση των γραμμών είναι πολλαπλάσια της αντίστασής τους, οπότε
$\Delta V \approx XQ/V_2$ και $\delta \approx XP/V_2^2$: η διαφορά των μέτρων των τάσεων
καθορίζεται από τη ροή άεργης ισχύος και η διαφορά γωνιών από τη ροή ενεργής ισχύος. Για
$R = 0$ η \eqref{eq:v1phasor}, αναλυμένη σε συνιστώσες, δίνει τις ακριβείς σχέσεις

\begin{equation}
P = \frac{V_1 V_2}{X}\,\sin\delta, \qquad
Q = \frac{V_1 V_2 \cos\delta - V_2^2}{X}
\label{eq:pq-line}
\end{equation}

όπου $Q$ η άεργη ισχύς που φθάνει στον ζυγό 2. Στα δίκτυα διανομής ο λόγος $X/R$ είναι
μικρότερος και η ενεργή ισχύς επηρεάζει σημαντικά και το μέτρο της τάσης· το θέμα εξετάζεται
στο Κεφάλαιο 5.

\begin{figure}[!htbp]
\centering
\includegraphics[width=\linewidth]{figures/v-q-phasor.svg}
\caption{Γραμμή δύο ζυγών με αντίδραση σειράς $X$ (αριστερά) και διανυσματικό διάγραμμα των
τάσεων (δεξιά) όταν η γραμμή τροφοδοτεί επαγωγικό φορτίο στον ζυγό 2, οπότε το ρεύμα
καθυστερεί της τάσης $V_2$. Η πτώση τάσης $jX\tilde{I}$
αναλύεται σε μια συνιστώσα $XQ/V_2$ σε φάση με την $V_2$, η οποία καθορίζει κυρίως τη
διαφορά των μέτρων, και σε μια κάθετη συνιστώσα $XP/V_2$, η οποία καθορίζει τη διαφορά
γωνιών $\delta$.}
\label{fig:vq}
\end{figure}

Από τις σχέσεις αυτές προκύπτουν τρία βασικά χαρακτηριστικά της ρύθμισης τάσης:

\begin{itemize}
\item Η άεργη ισχύς ρέει από τον ζυγό με τη μεγαλύτερη προς τον ζυγό με τη μικρότερη τάση,
      όπως η ενεργή ισχύς ρέει από τον ζυγό με τη μεγαλύτερη προς τον ζυγό με τη μικρότερη
      γωνία.
\item Η μεταφορά άεργης ισχύος σε μεγάλες αποστάσεις δεν είναι αποδοτική. Απαιτεί μεγάλες
      διαφορές τάσης και αυξάνει το ρεύμα, άρα τις απώλειες ενεργής ισχύος $RI^2$ και την
      άεργη ισχύ $XI^2$ που καταναλώνει η ίδια η γραμμή. Η άεργη ισχύς πρέπει επομένως να
      παράγεται κοντά στο σημείο όπου καταναλώνεται.
\item Η τάση είναι τοπικό μέγεθος. Ενώ η συχνότητα ρυθμίζεται από οποιαδήποτε μονάδα της
      συγχρονισμένης περιοχής, η τάση ενός ζυγού ρυθμίζεται αποτελεσματικά μόνο από μέσα
      που βρίσκονται ηλεκτρικά κοντά του.
\end{itemize}

Η επίδραση μιας μεταβολής $\Delta Q$ της άεργης ισχύος που εγχέεται σε έναν ζυγό εκτιμάται με
το ισοδύναμο Thévenin του δικτύου, όπως φαίνεται από τον ζυγό: μια πηγή τάσης πίσω από
αντίδραση $X_{\mathrm{th}}$. Από την \eqref{eq:dv} για $R = 0$ προκύπτει
$\Delta V \approx X_{\mathrm{th}}\,\Delta Q / V$ και, με την ισχύ βραχυκύκλωσης
$S_{\mathrm{sc}} = V^2 / X_{\mathrm{th}}$ του ζυγού,

\begin{equation}
\frac{\Delta V}{V} \approx \frac{\Delta Q}{S_{\mathrm{sc}}}
\label{eq:dv-ssc}
\end{equation}

Η ευαισθησία της τάσης στην άεργη ισχύ είναι επομένως αντιστρόφως ανάλογη της ισχύος
βραχυκύκλωσης. Σε ζυγό με μεγάλη ισχύ βραχυκύκλωσης, δηλαδή ηλεκτρικά κοντά σε πολλές
σύγχρονες μηχανές, η τάση μεταβάλλεται ελάχιστα, ενώ σε ασθενή ζυγό μικρές μεταβολές της
άεργης ισχύος μεταβάλλουν σημαντικά την τάση. Η σχέση συνδέει τη ρύθμιση τάσης με το ρεύμα
βραχυκύκλωσης, το οποίο αποτελεί χωριστή υπηρεσία (Πίνακας~\ref{tab:nfas}).

Οι παραπάνω σχέσεις αγνοούν την εγκάρσια χωρητικότητα των γραμμών, η οποία στις μεγάλες
γραμμές και στα καλώδια δεν μπορεί να αγνοηθεί. Η γραμμή γίνεται τότε η ίδια πηγή ή καταναλωτής
άεργης ισχύος: \emph{παράγει} άεργη ισχύ στη χωρητικότητά της, ανάλογη του τετραγώνου της
τάσης, και \emph{καταναλώνει} άεργη ισχύ στην αυτεπαγωγή της, ανάλογη του τετραγώνου του
ρεύματος. Για γραμμή μήκους $\ell$, με αυτεπαγωγή $L$ και χωρητικότητα $C$ ανά μονάδα μήκους,
και για τάση και ρεύμα περίπου σταθερά κατά μήκος της, η καθαρή άεργη ισχύς που παράγει η
γραμμή είναι

\begin{equation}
Q_{\mathrm{line}} \approx \omega \ell \left( C V^2 - L I^2 \right)
\label{eq:qline}
\end{equation}

Σε μια ελαφρά φορτισμένη γραμμή υπερισχύει ο πρώτος όρος. Η άεργη ισχύς που παράγει η
χωρητικότητα ρέει προς τα άκρα της γραμμής, οπότε, σύμφωνα με την \eqref{eq:dv}, η τάση στο
εσωτερικό της γραμμής υπερβαίνει τις τάσεις των άκρων. Σε γραμμή ανοικτή στο ένα άκρο η τάση
αυξάνεται προς το ανοικτό άκρο και υπερβαίνει την τάση τροφοδότησης· η αύξηση αυτή ονομάζεται
\emph{φαινόμενο Ferranti} και είναι ο λόγος για τον οποίο οι μεγάλες γραμμές υπερυψηλής τάσης
εξοπλίζονται με πηνία αντιστάθμισης στα άκρα τους. Σε μια φορτισμένη γραμμή υπερισχύει ο
δεύτερος όρος: η γραμμή καταναλώνει άεργη ισχύ, η οποία πρέπει να καλυφθεί από τα άκρα της, και
η τάση μειώνεται κατά μήκος της. Τα καλώδια έχουν πολύ μεγαλύτερη χωρητικότητα από τις εναέριες
γραμμές και παράγουν καθαρή άεργη ισχύ ακόμη και στη μέγιστη φόρτισή τους, οπότε οι μεγάλες
καλωδιακές συνδέσεις, όπως οι υποβρύχιες διασυνδέσεις νησιών, απαιτούν πηνία αντιστάθμισης.

Για τον ίδιο λόγο οι ανάγκες ρύθμισης τάσης μεταβάλλονται κατά τη διάρκεια της ημέρας. Στις
ώρες αιχμής οι γραμμές καταναλώνουν άεργη ισχύ, οπότε οι τάσεις τείνουν να μειωθούν και
απαιτείται παραγωγή άεργης ισχύος. Στις ώρες χαμηλού φορτίου οι γραμμές παράγουν άεργη ισχύ,
οπότε οι τάσεις τείνουν να αυξηθούν και απαιτείται απορρόφηση άεργης ισχύος. Οι υψηλές τάσεις
είναι συχνότερες όταν η διανεμημένη παραγωγή στα δίκτυα διανομής καλύπτει μεγάλο μέρος της
ζήτησης, οπότε μειώνεται η ροή ισχύος στο ΣΜΗΕ, καθώς και σε δίκτυα με πολλά καλώδια.

\subsection{Μέσα ρύθμισης της τάσης}
\label{sec:q-sources}

Τα μέσα ρύθμισης της τάσης διακρίνονται σε τρεις κατηγορίες:

\begin{itemize}
\item \textbf{Πηγές άεργης ισχύος με συνεχή ρύθμιση}: οι σύγχρονες γεννήτριες, οι μονάδες που
      συνδέονται μέσω μετατροπέων, οι σύγχρονοι αντισταθμιστές και οι διατάξεις ευέλικτων
      συστημάτων μεταφοράς εναλλασσόμενου ρεύματος (Flexible AC Transmission Systems, FACTS)
      παράλληλης σύνδεσης.
\item \textbf{Στατικά μέσα αντιστάθμισης με διακριτή ρύθμιση}: οι πυκνωτές και τα πηνία
      αντιστάθμισης και οι πυκνωτές σειράς.
\item \textbf{Μέσα που ρυθμίζουν τις τάσεις χωρίς να παράγουν/απορροφούν άεργη ισχύ}: οι
      μετασχηματιστές με σύστημα αλλαγής τάσης υπό φορτίο (ΣΑΤΥΦ) και οι αλλαγές τοπολογίας.
\end{itemize}

Τα μέσα αυτά διαφέρουν ως προς την ταχύτητα απόκρισης, τον τρόπο ρύθμισης και την εξάρτηση
της άεργης ισχύος από την τάση. Τα χαρακτηριστικά τους συνοψίζονται στον
Πίνακα~\ref{tab:v-means} στο τέλος της ενότητας.

\subsubsection{Σύγχρονες γεννήτριες}
\label{sec:sg}

Ο αυτόματος ρυθμιστής τάσης (Automatic Voltage Regulator, AVR) μετρά την τάση στους
ακροδέκτες της γεννήτριας, τη συγκρίνει με την τιμή αναφοράς και μεταβάλλει μέσω του διεγέρτη
την τάση διέγερσης, ώστε η τάση να παραμένει στην τιμή αναφοράς. Η γεννήτρια παράγει ή
απορροφά άεργη ισχύ όταν βρίσκεται σε κατάσταση υπερδιέγερσης ή υποδιέγερσης αντίστοιχα.
Την τιμή αναφοράς ορίζει ο χειριστής του σταθμού κατ' εντολή του ΔΣΜΗΕ ή, όπου
εφαρμόζεται, η δευτερεύουσα
ρύθμιση τάσης (Ενότητα «\nameref{sec:v-hier}»). Δύο συμπληρωματικές λειτουργίες του AVR
αφορούν άμεσα το δίκτυο:

\begin{itemize}
\item Η \emph{αντιστάθμιση άεργου ρεύματος} (reactive current compensation) εισάγει μικρό
      στατισμό στη ρύθμιση: η τάση αναφοράς μειώνεται ελαφρά όσο αυξάνεται η άεργη ισχύς της
      γεννήτριας. Έτσι οι μονάδες που συνδέονται στον ίδιο ζυγό μοιράζονται σταθερά την άεργη
      ισχύ, αντί να επιχειρεί καθεμία να επιβάλει τη δική της τάση.
\item Η \emph{αντιστάθμιση πτώσης τάσης} (line drop compensation) μεταφέρει το σημείο
      ρύθμισης πέρα από τους ακροδέκτες, συνήθως στην πλευρά υψηλής τάσης του μετασχηματιστή
      ανύψωσης, ώστε η γεννήτρια να ρυθμίζει αποτελεσματικότερα την τάση του δικτύου.
\end{itemize}

Και οι δύο λειτουργίες υλοποιούνται στη μέτρηση του AVR (Σχήμα~\ref{fig:avr}): αντί για την
τάση ακροδεκτών $V_t$, ο AVR ρυθμίζει την τάση $V_c = |\tilde{V}_t + jX_c\tilde{I}_t|$, όπου
$\tilde{I}_t$ το ρεύμα της γεννήτριας και $X_c$ η αντίδραση αντιστάθμισης. Με $X_c > 0$ η τάση
ακροδεκτών μειώνεται όσο αυξάνεται η παραγόμενη άεργη ισχύς, δηλαδή εισάγεται στατισμός
(αντιστάθμιση άεργου ρεύματος). Με $X_c < 0$, σε απόλυτη τιμή μικρότερη από την αντίδραση του
μετασχηματιστή ανύψωσης, η τάση ακροδεκτών αυξάνεται με την άεργη ισχύ, ώστε να
αντισταθμίζεται μέρος της πτώσης τάσης στον μετασχηματιστή (αντιστάθμιση πτώσης τάσης).

\begin{figure}[!htbp]
\centering
\includegraphics[width=\linewidth]{figures/avr-functions.svg}
\caption{Σύστημα διέγερσης σύγχρονης γεννήτριας με την αντιστάθμιση και τους περιοριστές
(αριστερά) και τάση ακροδεκτών συναρτήσει της παραγόμενης άεργης ισχύος για τα δύο πρόσημα
της αντίδρασης αντιστάθμισης $X_c$ (δεξιά). Η διακεκομμένη γραμμή αντιστοιχεί σε ρύθμιση της
τάσης ακροδεκτών χωρίς αντιστάθμιση.}
\label{fig:avr}
\end{figure}

Τα όρια λειτουργίας της γεννήτριας ορίζουν την \emph{περιοχή λειτουργίας σύγχρονης γεννήτριας}
(capability curve) του Σχήματος~\ref{fig:capability}. Για γεννήτρια κυλινδρικού δρομέα με
σύγχρονη αντίδραση $X_d$ και τάση ακροδεκτών $V$, τα δύο θερμικά όρια είναι κύκλοι στο
επίπεδο $P$--$Q$:

\begin{equation}
P^2 + Q^2 \le \left(V I_{\max}\right)^2, \qquad
P^2 + \left(Q + \frac{V^2}{X_d}\right)^2 \le \left(\frac{V E_{\max}}{X_d}\right)^2
\label{eq:capability}
\end{equation}

Η πρώτη σχέση είναι το \emph{όριο ρεύματος στάτη}, λόγω της θέρμανσης του τυλίγματος του
στάτη. Η δεύτερη είναι το \emph{όριο ρεύματος διέγερσης}, λόγω της θέρμανσης του τυλίγματος
του δρομέα, όπου $E_{\max}$ η εσωτερική τάση που αντιστοιχεί στο μέγιστο ρεύμα διέγερσης· το
όριο αυτό περιορίζει την παραγωγή άεργης ισχύος. Στην περιοχή υποδιέγερσης η απορρόφηση άεργης
ισχύος περιορίζεται από το \emph{όριο υποδιέγερσης}, λόγω ευστάθειας και θέρμανσης των άκρων
του πυρήνα του στάτη, ενώ η ενεργή ισχύς περιορίζεται από την κινητήρια μηχανή. Και τα δύο
θερμικά όρια εξαρτώνται από την τάση: σε χαμηλή τάση η μέγιστη άεργη ισχύς μειώνεται, δηλαδή
ακριβώς όταν χρειάζεται περισσότερο. Η διαφορά της μέγιστης άεργης ισχύος από την τρέχουσα
ονομάζεται \emph{εφεδρεία άεργης ισχύος} της μονάδας.

Ο Κώδικας Διαχείρισης του ΕΣΜΗΕ ορίζει την ελάχιστη ικανότητα άεργης ισχύος των υφιστάμενων
θερμικών και υδροηλεκτρικών μονάδων σε συνάρτηση με την τάση του δικτύου. Στη μέγιστη συνεχή
φόρτιση η μονάδα πρέπει να μπορεί να λειτουργεί από συντελεστή ισχύος 0,93 σε υποδιέγερση έως
0,85 σε υπερδιέγερση, για τάσεις 0,933--1,087 p.u. στο δίκτυο των 150 kV και 0,90--1,05 p.u.
στο δίκτυο των 400 kV (Σχήμα~\ref{fig:capability}).

\begin{figure}[!htbp]
\centering
\includegraphics[width=0.52\linewidth]{figures/capability.svg}
\caption{Περιοχή λειτουργίας σύγχρονης γεννήτριας (capability curve) ως
προς την ονομαστική φαινόμενη ισχύ και για τάση 1 p.u. Η περιοχή είναι ασύμμετρη, με μεγαλύτερη
ικανότητα παραγωγής από ό,τι απορρόφησης άεργης ισχύος· η κουκκίδα είναι το ονομαστικό σημείο
λειτουργίας, με συντελεστή ισχύος 0,8. Τα λευκά σημεία είναι οι απαιτήσεις του Κώδικα
Διαχείρισης του ΕΣΜΗΕ στη μέγιστη συνεχή φόρτιση, η οποία συνήθως συμπίπτει με το όριο του
στροβίλου.}
\label{fig:capability}
\end{figure}

Τα όρια της διέγερσης επιβάλλονται από περιοριστές του συστήματος διέγερσης
(Σχήμα~\ref{fig:avr}). Ο περιοριστής
υπερδιέγερσης (Over-Excitation Limiter, OEL) επιτρέπει υπερφόρτιση της διέγερσης για λίγα
δευτερόλεπτα, ώστε η γεννήτρια να στηρίζει την τάση κατά τις διαταραχές, και στη συνέχεια
επαναφέρει το ρεύμα διέγερσης στην επιτρεπόμενη τιμή με χαρακτηριστική αντίστροφου χρόνου.
Από τη στιγμή αυτή η γεννήτρια παύει να ρυθμίζει την τάση και συμπεριφέρεται ως πηγή σταθερής
άεργης ισχύος, γεγονός που συμβάλλει στην αστάθεια τάσης (Ενότητα
«\nameref{sec:v-stability}»). Ο περιοριστής υποδιέγερσης (Under-Excitation Limiter, UEL)
εμποδίζει αντίστοιχα τη λειτουργία πέρα από το όριο υποδιέγερσης. Τέλος, η άεργη ισχύς της
γεννήτριας φθάνει στο δίκτυο μέσω του μετασχηματιστή ανύψωσης, ο οποίος καταναλώνει μέρος
της, $X_T I^2$, όπου $X_T$ η αντίδρασή του.

\subsubsection{Μονάδες που συνδέονται μέσω μετατροπέων}
\label{sec:conv}

Τα αιολικά και φωτοβολταϊκά πάρκα, οι σταθμοί αποθήκευσης με συσσωρευτές και οι σταθμοί
συνεχούς ρεύματος υψηλής τάσης (High Voltage Direct Current, HVDC) με μετατροπείς πηγής τάσης
ρυθμίζουν την άεργη ισχύ τους ανεξάρτητα από την ενεργή, εντός του ορίου ρεύματος του
μετατροπέα:

\begin{equation}
P^2 + Q^2 \le \left(V I_{\max}\right)^2
\label{eq:conv-limit}
\end{equation}

Σε αντίθεση με τη σύγχρονη γεννήτρια, η περιοχή λειτουργίας είναι συμμετρική ως προς την
παραγωγή και την απορρόφηση άεργης ισχύος και συρρικνώνεται ανάλογα με την τάση
(Σχήμα~\ref{fig:capability-conv}). Ένας μετατροπέας με ονομαστική φαινόμενη ισχύ ίση με τη μέγιστη
ενεργή ισχύ της μονάδας δεν διαθέτει περιθώριο άεργης ισχύος όταν η μονάδα αποδίδει τη
μέγιστη ενεργή ισχύ της· γι' αυτό οι μετατροπείς υπερδιαστασιολογούνται. Αντίθετα, όταν η
ενεργή ισχύς είναι μηδενική, για παράδειγμα σε φωτοβολταϊκό σταθμό τη νύχτα, ο μετατροπέας
μπορεί να διαθέσει όλη την ονομαστική του ισχύ ως άεργη, εφόσον παραμένει συνδεδεμένος,
λειτουργώντας ουσιαστικά ως στατικός σύγχρονος αντισταθμιστής (Static Synchronous Compensator,
STATCOM), διάταξη FACTS που εξετάζεται στην Ενότητα «\nameref{sec:facts}».

\begin{figure}[!htbp]
\centering
\includegraphics[width=0.52\linewidth]{figures/capability-conv.svg}
\caption{Περιοχή λειτουργίας μονάδας που συνδέεται μέσω μετατροπέα, σε ανά μονάδα τιμές ως
προς την ονομαστική φαινόμενη ισχύ του μετατροπέα, για τάση 1 p.u. (συνεχής γραμμή) και
0,9 p.u. (διακεκομμένη). Η περιοχή είναι συμμετρική και περιορίζεται από το ρεύμα
του μετατροπέα και από τη μέγιστη ενεργή ισχύ $P_{\max}$ της μονάδας, οπότε συρρικνώνεται
με την τάση.}
\label{fig:capability-conv}
\end{figure}

Η άεργη ισχύς ρυθμίζεται με έναν από τους εξής τρόπους: σε σταθερή τιμή που ορίζει ο
διαχειριστής, με σταθερό συντελεστή ισχύος, οπότε μεταβάλλεται με την ενεργή ισχύ, ή ως
συνάρτηση της τάσης με στατισμό,

\begin{equation}
Q = \frac{U_{\mathrm{ref}} - U}{s\,U_n}\,Q_{\max}, \qquad -Q_{\max} \le Q \le Q_{\max}
\label{eq:qu}
\end{equation}

όπου $U$ η τάση στο σημείο σύνδεσης, $U_{\mathrm{ref}}$ η τάση αναφοράς και $s$ ο στατισμός,
δηλαδή η απόκλιση της τάσης, ως ποσοστό της ονομαστικής $U_n$, για την οποία
χρησιμοποιείται όλο το εύρος άεργης ισχύος (Σχήμα~\ref{fig:qu-droop}). Μόνο ο τρίτος τρόπος αποκρίνεται αυτόματα στις
μεταβολές της τάσης και συμβάλλει στην πρωτεύουσα ρύθμιση. Ο RfG, ο οποίος ονομάζει τις
μονάδες ηλεκτροπαραγωγής που συνδέονται μέσω μετατροπέων \emph{μονάδες πάρκου ισχύος} (power
park modules), ορίζει το πλαίσιο της απαιτούμενης ικανότητας άεργης ισχύος και των τρόπων
ρύθμισης, ενώ τις συγκεκριμένες τιμές καθορίζει ο αρμόδιος διαχειριστής. Η συμμετοχή των μονάδων των δικτύων διανομής στη ρύθμιση
της τάσης εξετάζεται στο Κεφάλαιο 5.

\begin{figure}[!htbp]
\centering
\includegraphics[width=0.52\linewidth]{figures/qu-droop.svg}
\caption{Χαρακτηριστική άεργης ισχύος-τάσης με στατισμό σύμφωνα με την \eqref{eq:qu}, για
$U_{\mathrm{ref}}$ = 1 p.u. και $s$ = 4\%. Η άεργη ισχύς μεταβάλλεται γραμμικά με την
απόκλιση της τάσης από την τάση αναφοράς και φθάνει το όριο $\pm Q_{\max}$ όταν η απόκλιση
γίνει ίση με $s\,U_n$.}
\label{fig:qu-droop}
\end{figure}

\subsubsection{Σύγχρονοι αντισταθμιστές}

Ο σύγχρονος αντισταθμιστής (synchronous condenser) είναι σύγχρονη μηχανή χωρίς κινητήρια
μηχανή, η οποία λειτουργεί ως σύγχρονος κινητήρας εν κενώ και ρυθμίζει την άεργη ισχύ της
μέσω της διέγερσης, με AVR όπως η γεννήτρια. Απορροφά από το δίκτυο μόνο την ενεργή ισχύ των
απωλειών της, ενώ η ικανότητα απορρόφησης άεργης ισχύος είναι συνήθως μικρότερη από την
ικανότητα παραγωγής. Εκτός από άεργη ισχύ, συνεισφέρει ρεύμα βραχυκύκλωσης και αδράνεια,
όπως κάθε σύγχρονη μηχανή. Γι' αυτό η χρήση του επεκτείνεται σε συστήματα με υψηλή διείσδυση
μονάδων με μετατροπείς, ενώ σε ορισμένες περιπτώσεις γεννήτριες σταθμών που αποσύρονται
μετατρέπονται σε σύγχρονους αντισταθμιστές.

\subsubsection{Πυκνωτές και πηνία αντιστάθμισης}

Οι πυκνωτές και τα πηνία αντιστάθμισης (shunt capacitors, shunt reactors) εγκαθίστανται σε
ζυγούς υποσταθμών ή στο τριτεύον τύλιγμα μετασχηματιστών και συνδέονται μέσω διακοπτών ισχύος
σε διακριτά βήματα. Η άεργη ισχύς τους είναι

\begin{equation}
Q = B V^2
\label{eq:shunt}
\end{equation}

όπου $B$ η επιδεκτικότητα της συστοιχίας, θετική για πυκνωτή και αρνητική για πηνίο. Η
μεταβολή της τάσης που προκαλεί η ζεύξη ενός βήματος εκτιμάται από την \eqref{eq:dv-ssc}. Οι
πυκνωτές εγκαθίστανται κυρίως κοντά στα κέντρα κατανάλωσης και συνδέονται στις ώρες αιχμής,
ενώ τα πηνία εγκαθίστανται στα άκρα μεγάλων γραμμών και καλωδίων και συνδέονται στις ώρες
χαμηλού φορτίου· τα πηνία των μεγάλων γραμμών είναι συχνά μόνιμα συνδεδεμένα. Σύμφωνα με την
\eqref{eq:shunt}, η άεργη ισχύς ενός πυκνωτή μειώνεται με το τετράγωνο της τάσης, δηλαδή
ακριβώς όταν η στήριξη της τάσης είναι περισσότερο αναγκαία. Επιπλέον, ο αριθμός των χειρισμών
είναι περιορισμένος, λόγω της καταπόνησης των διακοπτών και των μεταβατικών φαινομένων της
ζεύξης, οπότε τα μέσα αυτά καλύπτουν τις αργές μεταβολές των αναγκών και όχι τις ταχείες.

\subsubsection{Πυκνωτές σειράς}

Ο πυκνωτής σειράς αντισταθμίζει μέρος της επαγωγικής αντίδρασης $X_L$ μιας μεγάλης γραμμής,
με βαθμό αντιστάθμισης $k_{\mathrm{se}} = X_C / X_L$ συνήθως 25--70\%. Κύριος σκοπός του
είναι η αύξηση της μεταφορικής ικανότητας της γραμμής. Σύμφωνα με
την \eqref{eq:pq-line}, για δεδομένη διαφορά γωνιών η ενεργή ισχύς αυξάνεται κατά τον
παράγοντα $1/(1 - k_{\mathrm{se}})$· για δεδομένη ισχύ, αντίστοιχα, μειώνεται η διαφορά
γωνιών και αυξάνεται το περιθώριο ευστάθειας της μεταφοράς. Σε παράλληλες διαδρομές ο
πυκνωτής μεταβάλλει επίσης τον καταμερισμό της ροής ισχύος, υπέρ της γραμμής στην οποία
εγκαθίσταται. Η στήριξη της τάσης είναι δευτερεύον όφελος: μειώνεται η άεργη ισχύς που
καταναλώνει η γραμμή, ενώ η άεργη ισχύς που παράγει ο πυκνωτής, $X_C I^2$, αυξάνεται με το
τετράγωνο του ρεύματος, δηλαδή με τη φόρτιση της γραμμής. Η αντιστάθμιση προσαρμόζεται
επομένως αυτόματα στις ανάγκες.

\subsubsection{Μετασχηματιστές με σύστημα αλλαγής τάσης υπό φορτίο}
\label{sec:oltc}

Το σύστημα αλλαγής τάσης υπό φορτίο (ΣΑΤΥΦ· On-Load Tap Changer, OLTC) μεταβάλλει τον λόγο
μετασχηματισμού $n$ χωρίς διακοπή της τροφοδότησης, σε βήματα της τάξης του 1--1,5\% και
συνήθως σε εύρος $\pm$10--15\%. Αν αμεληθεί η αντίδραση του μετασχηματιστή, η τάση $V_2$ της
πλευράς που ρυθμίζεται και η τάση $V_1$ της άλλης πλευράς συνδέονται με τη σχέση $V_2 = n V_1$,
με τις τάσεις και τον λόγο $n$ σε ανά μονάδα τιμές, οπότε η αλλαγή λήψης μεταβάλλει άμεσα την
$V_2$.

Το ΣΑΤΥΦ δεν παράγει άεργη ισχύ· ανακατανέμει τις τάσεις και τις ροές άεργης ισχύος
μεταξύ των δύο πλευρών. Όταν αυξάνει την τάση της πλευράς του φορτίου, αυξάνονται η ενεργή και
η άεργη ισχύς του φορτίου, που εξαρτώνται από την τάση. Αυξάνεται επομένως η ροή ενεργής και άεργης ισχύος από το δίκτυο μεταφοράς οδηγώντας έτσι σε πτώση τάσης στο δίκτυο αυτό. Η αύξηση της τάσης του φορτίου γίνεται δηλαδή
σε βάρος της τάσης του δικτύου μεταφοράς. Στους αυτομετασχηματιστές 400/150 kV το ΣΑΤΥΦ ρυθμίζει την τάση
του δικτύου 150 kV και την ανταλλαγή άεργης ισχύος μεταξύ των δύο επιπέδων τάσης, ενώ στους
μετασχηματιστές 150/20 kV ρυθμίζει την τάση των δικτύων διανομής.

Την ίδια εξάρτηση της ζήτησης από την τάση αξιοποιεί αντίστροφα η \emph{μείωση της τάσης για
εξοικονόμηση ενέργειας} (Conservation Voltage Reduction, CVR): τα ΣΑΤΥΦ των μετασχηματιστών
150/20 kV ρυθμίζουν την τάση των δικτύων διανομής στο κατώτερο τμήμα του επιτρεπόμενου εύρους,
ώστε να μειώνονται η ζήτηση και η κατανάλωση ενέργειας, συνήθως κατά 0,4--1\% για κάθε 1\%
μείωσης της τάσης. Η μείωση της τάσης χρησιμοποιείται επίσης ως έκτακτο μέτρο
περιορισμού της ζήτησης σε συνθήκες ελλείμματος παραγωγής.

\subsubsection{Ευέλικτα συστήματα μεταφοράς εναλλασσόμενου ρεύματος (FACTS)}
\label{sec:facts}

Πέρα από τα συμβατικά μέσα, ο ΔΣΜΗΕ ελέγχει άμεσα τις ροές ισχύος και τις τάσεις του δικτύου
μεταφοράς με διατάξεις \textbf{FACTS}, δηλαδή διατάξεις ηλεκτρονικών ισχύος που επιτρέπουν τον
γρήγορο και συνεχή έλεγχο των μεγεθών που καθορίζουν τη ροή ισχύος στο δίκτυο. Για τη ρύθμιση της τάσης χρησιμοποιούνται κυρίως οι διατάξεις παράλληλης σύνδεσης,
ενώ οι διατάξεις σειράς ελέγχουν κυρίως τις ροές ενεργής ισχύος.

Σύμφωνα με την \eqref{eq:pq-line}, η ενεργή ισχύς που μεταφέρεται μεταξύ δύο ζυγών μέσω
γραμμής χωρίς απώλειες εξαρτάται από τα μέτρα των τάσεων $V_1, V_2$, την αντίδραση σειράς
$X$ και τη διαφορά γωνιών $\theta_1 - \theta_2$. Σε
συμβατικό δίκτυο τα μεγέθη αυτά δεν ελέγχονται άμεσα, οπότε η κατανομή των ροών καθορίζεται
από τις σύνθετες αντιστάσεις των κλάδων και όχι από τις απαιτήσεις λειτουργίας. Οι διατάξεις
FACTS καθιστούν ελεγχόμενο ένα ή περισσότερα από τα μεγέθη αυτά και κατατάσσονται ανάλογα με
το μέγεθος που επηρεάζουν (Σχήμα~\ref{fig:facts}):

\begin{itemize}
\item \textbf{Παράλληλης σύνδεσης} (έλεγχος του $V$): ο στατικός αντισταθμιστής άεργης ισχύος
      (Static Var Compensator, SVC) και ο στατικός σύγχρονος αντισταθμιστής (Static
      Synchronous Compensator, STATCOM). Ρυθμίζουν την τάση του ζυγού μέσω έγχυσης ή
      απορρόφησης άεργης ισχύος.
\item \textbf{Σειριακής σύνδεσης} (έλεγχος του $X$): ο πυκνωτής σειράς ελεγχόμενος με
      θυρίστορ (Thyristor Controlled Series Capacitor, TCSC) και ο στατικός σύγχρονος
      αντισταθμιστής σειράς (Static Synchronous Series Compensator, SSSC). Μεταβάλλουν την
      ισοδύναμη αντίδραση σειράς της γραμμής, ανακατανέμουν τις ροές ισχύος και αυξάνουν τη
      μεταφορική ικανότητα.
\item \textbf{Συνδυασμένης σύνδεσης}, σε σειρά και παράλληλα (έλεγχος και της γωνίας): ο
      ενοποιημένος ελεγκτής ροής ισχύος (Unified Power Flow Controller, UPFC) ελέγχει
      ταυτόχρονα και ανεξάρτητα την τάση και τις ροές ενεργής και άεργης ισχύος. Παρέχει τη
      μεγαλύτερη ευελιξία ελέγχου αλλά με υψηλότερο κόστος.
\end{itemize}

\begin{figure}[!htbp]
\centering
\includegraphics[width=\linewidth]{figures/facts.svg}
\caption{Οι τρεις κατηγορίες διατάξεων FACTS σε γραμμή δύο ζυγών. Κάθε κατηγορία επεμβαίνει
σε διαφορετικό μέγεθος της σχέσης μεταφοράς ισχύος, το οποίο σημειώνεται με έντονο χρώμα: η
παράλληλη στην τάση, η σειριακή στην αντίδραση, η συνδυασμένη και στη γωνία.}
\label{fig:facts}
\end{figure}

\textbf{Τοπολογίες διατάξεων.} Ανάλογα με την αρχή λειτουργίας τους, οι διατάξεις FACTS διακρίνονται σε διατάξεις
\emph{μεταβλητής σύνθετης αντίστασης}, με θυρίστορ, και σε διατάξεις
\emph{ελεγχόμενης πηγής τάσης}, με μετατροπείς πηγής τάσης. Τα μονογραμμικά διαγράμματα
δίνονται στο Σχήμα~\ref{fig:facts-devices}.

\begin{itemize}
\item Ο \textbf{SVC} αποτελείται από πηνίο ελεγχόμενο με θυρίστορ (Thyristor Controlled
      Reactor, TCR) και πυκνωτές μεταγόμενους με θυρίστορ (Thyristor Switched Capacitor,
      TSC), σε παράλληλη σύνδεση. Η γωνία έναυσης του TCR ρυθμίζει συνεχώς την ισοδύναμη
      επαγωγική επιδεκτικότητα, ενώ οι TSC προσθέτουν χωρητική επιδεκτικότητα σε διακριτά
      βήματα. Η διάταξη συμπεριφέρεται ως \emph{μεταβλητή επιδεκτικότητα} $B$ (susceptance).
\item Ο \textbf{STATCOM} είναι μετατροπέας πηγής τάσης (Voltage Source Converter, VSC) με
      πυκνωτή στην πλευρά συνεχούς ρεύματος (ΣΡ). Παράγει εναλλασσόμενη τάση $V_\mu$
      ελεγχόμενου μέτρου, η οποία εφαρμόζεται στο δίκτυο μέσω της αντίδρασης σύζευξης: για
      $V_\mu > V$ ο STATCOM εγχέει άεργη ισχύ, ενώ για $V_\mu < V$ απορροφά.
\item Ο \textbf{TCSC} είναι πυκνωτής σειράς με παράλληλο TCR, οπότε η ισοδύναμη αντίδραση
      σειράς της γραμμής μεταβάλλεται συνεχώς. Ο \textbf{SSSC} εγχέει τάση σε σειρά μέσω
      μετασχηματιστή σύζευξης σειράς· η εγχεόμενη τάση είναι ανεξάρτητη του ρεύματος της
      γραμμής, ενώ στον TCSC η τάση στα άκρα του πυκνωτή είναι ανάλογη του ρεύματος.
\item Ο \textbf{UPFC} αποτελείται από έναν παράλληλο και έναν σειριακό VSC με \emph{κοινό
      ζυγό ΣΡ} (DC link). Ο κοινός ζυγός επιτρέπει την ανταλλαγή \emph{ενεργής} ισχύος
      μεταξύ των δύο μετατροπέων, οπότε ο σειριακός μετατροπέας εγχέει τάση οποιουδήποτε
      μέτρου και γωνίας εντός των ονομαστικών του ορίων και ελέγχει ανεξάρτητα τις ροές
      ενεργής και άεργης ισχύος. Οι δύο μετατροπείς βρίσκονται στον ίδιο υποσταθμό, στο άκρο
      της γραμμής που ελέγχεται: ο παράλληλος συνδέεται στον ζυγό και ο σειριακός
      παρεμβάλλεται στη γραμμή μέσω του μετασχηματιστή σειράς, αμέσως μετά τον ζυγό.
\end{itemize}

Το πλεονέκτημα του UPFC φαίνεται στη γραμμή της \eqref{eq:pq-line}. Αν ο UPFC βρίσκεται στο
άκρο 1, η τάση $V_\sigma$ του σειριακού μετατροπέα προστίθεται στην $V_1$ και, για δεδομένα μέτρα
των τάσεων των δύο ζυγών και δεδομένη διαφορά γωνιών $\delta$, μετατοπίζει τη μιγαδική ισχύ
που φθάνει στον ζυγό 2 κατά $V_2 V_\sigma/X$, σε διεύθυνση που καθορίζει η γωνία της $V_\sigma$. Το
σημείο λειτουργίας μπορεί επομένως να βρεθεί οπουδήποτε μέσα σε κύκλο του επιπέδου $P$--$Q$, με κέντρο το σημείο
λειτουργίας χωρίς UPFC και ακτίνα $V_2 V_\sigma^{\max}/X$ (Σχήμα~\ref{fig:upfc-pq}).
Αντίθετα, οι TCSC και SSSC έχουν μία μόνο μεταβλητή ελέγχου και μετακινούν το σημείο
λειτουργίας μόνο πάνω σε μία ευθεία που διέρχεται από την αρχή των αξόνων, οπότε η ροή άεργης
ισχύος καθορίζεται από τη ροή ενεργής ισχύος. Ο UPFC μπορεί έτσι, π.χ., να
διατηρεί την ενεργή ισχύ της γραμμής στην προγραμματισμένη τιμή και ταυτόχρονα να μηδενίζει
την άεργη ισχύ που φθάνει στον ζυγό 2, ώστε να περιορίζονται το ρεύμα και οι απώλειες της
γραμμής.

\begin{figure}[!htbp]
\centering
\includegraphics[width=\linewidth]{figures/facts-devices.svg}
\caption{Μονογραμμικά διαγράμματα των βασικών διατάξεων FACTS. Οι δύο πρώτες συνδέονται
παράλληλα, οι δύο επόμενες σε σειρά, ενώ ο UPFC συνδυάζει αμφότερους τους τρόπους μέσω
κοινού ζυγού ΣΡ.}
\label{fig:facts-devices}
\end{figure}

\begin{figure}[!htbp]
\centering
\includegraphics[width=0.52\linewidth]{figures/upfc-pq.svg}
\caption{Περιοχή λειτουργίας στον ζυγό 2 της γραμμής της \eqref{eq:pq-line}, για $V_1 = V_2$ =
1 p.u., $X$ = 0,5 p.u., $\delta$ = 20° και $V_\sigma^{\max}$ = 0,2 p.u. Οι $P$ και $Q$ είναι η
ενεργή και η άεργη ισχύς που φθάνουν στον ζυγό 2· για $Q > 0$ ο ζυγός 2 απορροφά άεργη ισχύ
από τη γραμμή. Όλες οι διατάξεις σειράς βρίσκονται στο άκρο 1. Ο UPFC μετακινεί το σημείο
λειτουργίας A οπουδήποτε μέσα στον κύκλο, π.χ. στο B, με την ίδια ενεργή
ισχύ και μηδενική άεργη ισχύ. Οι TCSC και SSSC το μετακινούν μόνο πάνω στην ευθεία που
διέρχεται από την αρχή των αξόνων· το έντονο τμήμα της αντιστοιχεί σε SSSC με την ίδια μέγιστη
τάση $V_\sigma^{\max}$.}
\label{fig:upfc-pq}
\end{figure}

\textbf{Χαρακτηριστικές $V$--$I$ των SVC και STATCOM.} Η διαφορά στην αρχή λειτουργίας αποτυπώνεται στις χαρακτηριστικές $V$--$I$ των δύο
διατάξεων παράλληλης σύνδεσης (Σχήμα~\ref{fig:facts-vi}).

Ο SVC συμπεριφέρεται ως μεταβλητή επιδεκτικότητα, οπότε $I = B\,V$ και $Q = B\,V^2$. Στο
όριο χωρητικής λειτουργίας το ρεύμα είναι ανάλογο της τάσης του ζυγού: στο 50\% της
ονομαστικής τάσης ο SVC αποδίδει το 50\% του ονομαστικού ρεύματος και το 25\% της ονομαστικής
άεργης ισχύος. Ο STATCOM, ως πηγή τάσης με έλεγχο ρεύματος, διατηρεί το ονομαστικό ρεύμα έως
πολύ χαμηλές τιμές τάσης, με περιορισμό μόνο από το μέγιστο επιτρεπόμενο ρεύμα των
ημιαγωγικών διακοπτών· η άεργη ισχύς του μειώνεται επομένως γραμμικά και όχι τετραγωνικά
με την τάση.

Κατά τη διάρκεια και αμέσως μετά από βραχυκύκλωμα, όταν η τάση είναι χαμηλή και η ανάγκη
στήριξης τάσης μέγιστη, ο STATCOM αποδίδει σημαντικά μεγαλύτερη άεργη ισχύ από SVC ίδιας
ονομαστικής ισχύος.

\begin{figure}[!htbp]
\centering
\includegraphics[width=\linewidth]{figures/facts-vi.svg}
\caption{Χαρακτηριστικά $V$--$I$ των δύο διατάξεων παράλληλης σύνδεσης. Η σκιασμένη περιοχή
είναι η περιοχή λειτουργίας και η έντονη γραμμή η χαρακτηριστική λειτουργίας με στατισμό. Τα
όρια του SVC είναι ευθείες σταθερής επιδεκτικότητας που διέρχονται από την αρχή, ενώ του
STATCOM είναι κατακόρυφες ευθείες σταθερού ρεύματος.}
\label{fig:facts-vi}
\end{figure}

Οι διατάξεις FACTS, μαζί με τα συστήματα HVDC, στα οποία η ροή ισχύος καθορίζεται απευθείας από τον έλεγχο των
μετατροπέων, συμπληρώνουν τα συμβατικά μέσα του ΣΜΗΕ: για τη ρύθμιση της τάσης
τους ρυθμιστές τάσης των γεννητριών, τα πηνία και τους πυκνωτές αντιστάθμισης και τα
ΣΑΤΥΦ των μετασχηματιστών, και για τη διαχείριση των συμφορήσεων την ανακατανομή
παραγωγής, τις αλλαγές τοπολογίας και τους μετασχηματιστές μετατόπισης φάσης. Τα αντίστοιχα μέσα στα
δίκτυα διανομής εξετάζονται στα Κεφάλαια 5 και 6.

\subsubsection{Αλλαγές τοπολογίας}

Σε συνθήκες πολύ χαμηλού φορτίου οι ΔΣΜΗΕ θέτουν ενίοτε εκτός λειτουργίας ελαφρά φορτισμένες
γραμμές υπερυψηλής τάσης, ώστε να περιοριστεί η άεργη ισχύς που παράγουν. Η αλλαγή αυτή της
τοπολογίας επιτρέπεται μόνο εφόσον εξακολουθεί να
ικανοποιείται το κριτήριο N-1 (Κεφάλαιο 1, Ενότητα «\nameref{sec:n1}»).

\subsubsection{Σύγκριση των μέσων ρύθμισης}

Ο Πίνακας~\ref{tab:v-means} συνοψίζει τα χαρακτηριστικά των μέσων. Τα ταχέα μέσα με συνεχή
ρύθμιση αντιμετωπίζουν τις διαταραχές, ενώ τα βραδέα μέσα με διακριτή ρύθμιση καλύπτουν τις
αργές μεταβολές των αναγκών, ώστε τα ταχέα να διατηρούν την εφεδρεία τους. Σημαντική είναι και
η εξάρτηση της άεργης ισχύος από την τάση στο όριο κάθε μέσου: όσο ασθενέστερη είναι η
εξάρτηση αυτή, τόσο αποτελεσματικότερο είναι το μέσο σε χαμηλές τάσεις.

\begin{table}[!htbp]
\centering
\begin{tabular}{llll}
\hline
Μέσο & Απόκριση & Ρύθμιση & Μέγιστη άεργη ισχύς \\
\hline
σύγχρονη γεννήτρια & $<$ 1 s & συνεχής & περιοχή λειτουργίας \\
σύγχρονος αντισταθμιστής & $<$ 1 s & συνεχής & περιοχή λειτουργίας \\
μετατροπέας μονάδας ΑΠΕ & ms & συνεχής & ανάλογη της $V$ \\
STATCOM & ms & συνεχής & ανάλογη της $V$ \\
SVC & ms & συνεχής & ανάλογη του $V^2$ \\
TCSC & ms & συνεχής & ανάλογη του $I^2$ \\
SSSC & ms & συνεχής & ανάλογη του $I$ \\
UPFC & ms & συνεχής & όπως STATCOM και SSSC \\
πυκνωτής, πηνίο αντιστάθμισης & s -- min & διακριτή & ανάλογη του $V^2$ \\
πυκνωτής σειράς & s -- min & διακριτή & ανάλογη του $I^2$ \\
ΣΑΤΥΦ & δεκάδες s & διακριτή & δεν παράγει \\
αλλαγή τοπολογίας & min & διακριτή & δεν παράγει \\
\hline
\end{tabular}
\caption{Χαρακτηριστικά των μέσων ρύθμισης της τάσης. Η τελευταία στήλη δίνει την εξάρτηση της
μέγιστης άεργης ισχύος από την τάση ή, για τα μέσα σειράς, από το ρεύμα $I$ της γραμμής.}
\label{tab:v-means}
\end{table}

\subsection{Ρύθμιση τάσης από τον προγραμματισμό στον πραγματικό χρόνο}
\label{sec:v-stages}

Όπως το ισοζύγιο ισχύος (Κεφάλαιο 1, Ενότητα «\nameref{sec:stages}»), έτσι και η ρύθμιση της
τάσης προετοιμάζεται σε πολλαπλούς χρονικούς ορίζοντες. Οι αποφάσεις κάθε σταδίου εξασφαλίζουν
ότι στο επόμενο υπάρχουν τα αναγκαία μέσα για τη αποτελεσματική ρύθμιση της τάσης: η ανάπτυξη του δικτύου εξασφαλίζει τον εξοπλισμό, ο
προγραμματισμός της λειτουργίας τη διαθεσιμότητά του και ο πραγματικός χρόνος τη χρήση του.

\begin{itemize}
\item \textbf{Ανάπτυξη του δικτύου} (έτη). Μελέτες ροής ισχύος και ευστάθειας τάσης για
      σενάρια ζήτησης και παραγωγής καθορίζουν τις θέσεις και τα μεγέθη νέων πηνίων, πυκνωτών,
      STATCOM και σύγχρονων αντισταθμιστών. Οι απαιτήσεις σύνδεσης εξασφαλίζουν την ικανότητα
      άεργης ισχύος των νέων μονάδων.
\item \textbf{Προγραμματισμός της λειτουργίας} (μήνες έως εβδομάδες). Τα προγράμματα
      συντήρησης γραμμών, μετασχηματιστών, μονάδων και μέσων αντιστάθμισης συντονίζονται,
      ώστε καμία περιοχή να μη μείνει χωρίς επαρκή στήριξη τάσης. Εντοπίζονται οι περίοδοι με
      κίνδυνο υψηλών τάσεων, όπως οι αργίες και οι νύχτες με χαμηλή ζήτηση, και οι περίοδοι με
      κίνδυνο χαμηλών τάσεων, όπως οι αιχμές του καλοκαιριού.
\item \textbf{Επόμενη ημέρα}. Ο ΔΣΜΗΕ ελέγχει την τεχνική εφικτότητα του προγράμματος που
      προκύπτει από τις αγορές: για κάθε αγοραία χρονική μονάδα της επόμενης ημέρας εκτελεί
      ροές ισχύος και ανάλυση απρόβλεπτων συμβάντων σε μοντέλο του δικτύου με τις προβλέψεις φορτίου
      και ΑΠΕ και ελέγχει τις τάσεις στις καταστάσεις N και N-1. Όπου απαιτείται, προβλέπει την
      ένταξη μονάδων σε περιοχές με ανεπαρκή στήριξη τάσης, καθορίζει τις τάσεις αναφοράς ή την
      άεργη ισχύ των μονάδων και καταρτίζει πρόγραμμα χειρισμών για πηνία, πυκνωτές και
      γραμμές.
\item \textbf{Ενδοημερήσιος ορίζοντας}. Οι έλεγχοι επαναλαμβάνονται με ενημερωμένες
      προβλέψεις φορτίου και ΑΠΕ και τα προγράμματα αναθεωρούνται.
\item \textbf{Πραγματικός χρόνος}. Το SCADA συλλέγει τις τάσεις, τις ροές άεργης ισχύος και τις
      θέσεις των ΣΑΤΥΦ, η εκτίμηση κατάστασης δίνει τις τάσεις όλων των ζυγών, η
      ανάλυση απρόβλεπτων συμβάντων ελέγχει τις τάσεις μετά από κάθε συμβάν του καταλόγου και
      η βέλτιστη ροή ισχύος προτείνει τιμές αναφοράς. Οι χειριστές εκδίδουν εντολές για τις
      τάσεις αναφοράς ή την άεργη ισχύ των μονάδων, συνδέουν ή αποσυνδέουν πηνία και πυκνωτές, μεταβάλλουν
      τις τιμές αναφοράς των ΣΑΤΥΦ και, όταν απαιτείται, αλλάζουν την τοπολογία.
      Παράλληλα λειτουργούν συνεχώς οι αυτόματοι ρυθμιστές των μονάδων και των διατάξεων
      αντιστάθμισης, οργανωμένοι στα επίπεδα της ιεραρχικής ρύθμισης τάσης (Ενότητα
      «\nameref{sec:v-hier}»), ενώ σε κατάσταση έκτακτης ανάγκης εφαρμόζονται τα μέτρα του
      σχεδίου άμυνας του συστήματος.
\item \textbf{Μετά τη λειτουργία}. Αναλύονται τα συμβάντα και παρακολουθείται η συμμόρφωση των
      μονάδων με τις απαιτήσεις άεργης ισχύος και τις εντολές του ΔΣΜΗΕ.
\end{itemize}

Σε αντίθεση με την εξισορρόπηση, η άεργη ισχύς δεν μεταφέρεται αποτελεσματικά σε μεγάλες
αποστάσεις (Ενότητα «\nameref{sec:v-q}»), οπότε κάθε απόφαση αφορά συγκεκριμένη περιοχή του
δικτύου και συγκεκριμένα μέσα σε αυτήν. Για τον λόγο αυτό τα εργαλεία κάθε σταδίου βασίζονται
σε πλήρες μοντέλο του δικτύου, με τα μέτρα των τάσεων και την άεργη
ισχύ, και όχι στο γραμμικό μοντέλο ΣΡ, το οποίο αγνοεί και τα δύο.

\subsection{Ιεραρχική ρύθμιση τάσης}
\label{sec:v-hier}

Η ρύθμιση της τάσης οργανώνεται ιεραρχικά σε τρία επίπεδα, τα οποία διαφέρουν ως προς τη
γεωγραφική έκταση και τη χρονική κλίμακα (Σχήμα~\ref{fig:v-hier} και
Πίνακας~\ref{tab:v-hier}). Το ανώτερο επίπεδο ορίζει τις τιμές αναφοράς του αμέσως
κατώτερου, και κάθε επίπεδο είναι σημαντικά βραδύτερο από το προηγούμενο, ώστε να
αντιλαμβάνεται τα ταχύτερα επίπεδα σε μόνιμη κατάσταση και να αποφεύγονται ανεπιθύμητες
αλληλεπιδράσεις. Η οργάνωση θυμίζει τη ρύθμιση φορτίου-συχνότητας
(\hyperref[kef-03-theoria]{Κεφάλαιο 3}), με τη διαφορά ότι η τάση είναι τοπικό μέγεθος, οπότε
τα επίπεδα της ρύθμισης ορίζονται ανά περιοχή και όχι για ολόκληρη τη συγχρονισμένη περιοχή.

\begin{figure}[!htbp]
\centering
\includegraphics[width=\linewidth]{figures/v-hier.svg}
\caption{Ιεραρχική ρύθμιση τάσης. Η τριτεύουσα ρύθμιση υπολογίζει στο κέντρο ελέγχου τις
τάσεις αναφοράς των ζυγών-πιλότων (pilot buses· Ενότητα «\nameref{sec:v-secondary}»), η
δευτερεύουσα ρύθμιση κάθε ζώνης διατηρεί την τάση του ζυγού-πιλότου μεταβάλλοντας ομοιόμορφα
την άεργη ισχύ των μονάδων της ζώνης και η πρωτεύουσα
ρύθμιση κάθε μονάδας διατηρεί την τάση στους ακροδέκτες της.}
\label{fig:v-hier}
\end{figure}

\begin{table}[!htbp]
\centering
\begin{tabular}{llll}
\hline
Επίπεδο & Έκταση & Χρονική κλίμακα & Κύρια μέσα \\
\hline
πρωτεύουσα   & τοπική  & ms -- s     & AVR, SVC, STATCOM, μετατροπείς \\
δευτερεύουσα & ζώνη    & min         & τιμές αναφοράς AVR, πυκνωτές, πηνία \\
τριτεύουσα   & σύστημα & 15 min -- h & βέλτιστη ροή ισχύος \\
\hline
\end{tabular}
\caption{Επίπεδα της ιεραρχικής ρύθμισης τάσης.}
\label{tab:v-hier}
\end{table}

\subsubsection{Πρωτεύουσα ρύθμιση}

Η πρωτεύουσα ρύθμιση είναι τοπική και αυτόματη. Την εκτελούν οι AVR των γεννητριών και των
σύγχρονων αντισταθμιστών, οι διατάξεις FACTS και οι μετατροπείς με ρύθμιση τάσης κατά την
\eqref{eq:qu}. Αποκρίνεται σε κλάσματα του δευτερολέπτου έως λίγα δευτερόλεπτα και διατηρεί
την τάση κάθε ρυθμιζόμενου ζυγού κοντά στην τιμή αναφοράς του, με μικρό στατισμό, ώστε
γειτονικά μέσα να μοιράζονται την άεργη ισχύ. Οι πρωτεύοντες ρυθμιστές δεν συντονίζονται μεταξύ
τους: καθένας ρυθμίζει τον δικό του ζυγό, ενώ η τάση των υπόλοιπων ζυγών προκύπτει από το
δίκτυο. Για τον λόγο αυτό, μετά από μια μεταβολή της ζήτησης η άεργη ισχύς κατανέμεται άνισα,
ανάλογα με την ηλεκτρική απόσταση κάθε μονάδας από τη μεταβολή, και οι πλησιέστερες μονάδες
μπορεί να φθάσουν στα όριά τους ενώ άλλες διαθέτουν ακόμη εφεδρεία.

\subsubsection{Δευτερεύουσα ρύθμιση}
\label{sec:v-secondary}

Η δευτερεύουσα ρύθμιση συντονίζει τις μονάδες κάθε περιοχής. Το σύστημα χωρίζεται σε
\emph{ζώνες} ρύθμισης τάσης, ηλεκτρικά ασθενώς συζευγμένες μεταξύ τους, και σε κάθε ζώνη
επιλέγεται ένας αντιπροσωπευτικός ζυγός, ο \emph{ζυγός-πιλότος} (pilot bus), η τάση του οποίου
μεταβάλλεται όπως οι τάσεις της ζώνης. Ο ρυθμιστής της ζώνης συγκρίνει την τάση $V_p$ του
ζυγού-πιλότου με την τιμή αναφοράς $V_p^{\mathrm{ref}}$ και υπολογίζει, με αναλογικό και
ολοκληρωτικό έλεγχο, ένα κοινό \emph{επίπεδο άεργης ισχύος} $q$ για όλες τις μονάδες της ζώνης:

\begin{equation}
q = K_P \left(V_p^{\mathrm{ref}} - V_p\right)
  + K_I \int \left(V_p^{\mathrm{ref}} - V_p\right) \mathrm{d}t, \qquad
Q_i^{\mathrm{ref}} = q\,Q_i^{\max}
\label{eq:svr}
\end{equation}

όπου $-1 \le q \le 1$ και $Q_i^{\max}$ η ικανότητα άεργης ισχύος της μονάδας $i$. Ένας τοπικός
βρόχος σε κάθε μονάδα μεταβάλλει την τιμή αναφοράς του AVR της, ώστε η άεργη ισχύς της να
ακολουθεί την $Q_i^{\mathrm{ref}}$. Όλες οι μονάδες της ζώνης χρησιμοποιούν έτσι το ίδιο ποσοστό
της ικανότητάς τους, με αποτέλεσμα ισόρροπη κατανομή των εφεδρειών άεργης ισχύος και καλύτερο
προφίλ τάσεων. Ο χρόνος απόκρισης είναι της τάξης του λεπτού, σημαντικά μεγαλύτερος από εκείνον
της πρωτεύουσας ρύθμισης. Αυτόματη δευτερεύουσα ρύθμιση εφαρμόζεται σε ορισμένα συστήματα,
όπως της Γαλλίας και της Ιταλίας. Σε πολλά συστήματα τον ρόλο της εκτελούν οι χειριστές του
κέντρου ελέγχου, με ζεύξη πυκνωτών και πηνίων, αλλαγή των τάσεων αναφοράς των μονάδων και
χειρισμούς των ΣΑΤΥΦ.

\subsubsection{Τριτεύουσα ρύθμιση}

Η τριτεύουσα ρύθμιση είναι κεντρική και βελτιστοποιεί τη λειτουργία ολόκληρου του συστήματος σε
χρονική κλίμακα δεκάδων λεπτών έως ωρών. Με τη βέλτιστη ροή ισχύος του EMS υπολογίζονται οι
μεταβλητές ελέγχου $\mathbf{u}$, δηλαδή οι τάσεις αναφοράς των μονάδων ή των ζυγών-πιλότων, οι
λήψεις των μετασχηματιστών και οι συστοιχίες πυκνωτών και πηνίων σε λειτουργία, από το πρόβλημα

\begin{equation}
\begin{aligned}
\min_{\mathbf{u}} \quad & P_{\mathrm{loss}}(\mathbf{x}, \mathbf{u}) \\
\text{υπό} \quad & \mathbf{g}(\mathbf{x}, \mathbf{u}) = \mathbf{0}, \qquad
V_i^{\min} \le V_i \le V_i^{\max}, \qquad
Q_{g,j}^{\min} \le Q_{g,j} \le Q_{g,j}^{\max}
\end{aligned}
\label{eq:orpf}
\end{equation}

όπου $\mathbf{x}$ οι τάσεις των ζυγών, $\mathbf{g}$ οι εξισώσεις ροής ισχύος, $P_{\mathrm{loss}}$
οι απώλειες ενεργής ισχύος και $Q_{g,j}$ η άεργη ισχύς της μονάδας $j$. Η αντικειμενική
συνάρτηση συμπληρώνεται συχνά με όρους που ευνοούν τη διατήρηση εφεδρείας άεργης ισχύος στις
γρήγορες πηγές, ενώ οι περιορισμοί μπορούν να επεκταθούν στις καταστάσεις μετά από κάθε
απρόβλεπτο συμβάν του καταλόγου, όπως στη βέλτιστη ροή ισχύος με περιορισμούς ασφάλειας
(Κεφάλαιο 1, Ενότητα «\nameref{sec:n1}»). Οι λήψεις και οι συστοιχίες είναι διακριτές
μεταβλητές, γεγονός που δυσχεραίνει την επίλυση. Τα αποτελέσματα αποτελούν τις τιμές αναφοράς
της δευτερεύουσας ρύθμισης ή, όπου αυτή δεν υπάρχει, οδηγίες προς τους χειριστές.

\subsubsection{Συντονισμός των επιπέδων και των μέσων}

Δύο αρχές συντονισμού είναι ιδιαίτερα σημαντικές. Πρώτον, η ταχεία εφεδρεία άεργης ισχύος των
γεννητριών και των διατάξεων FACTS πρέπει να διατηρείται για την αντιμετώπιση απρόβλεπτων
συμβάντων. Τα βραδέα μέσα, όπως οι πυκνωτές και τα πηνία, αναλαμβάνουν σταδιακά τη μόνιμη
κάλυψη των αναγκών, ώστε οι ταχείες πηγές να επανέρχονται κοντά στο μέσο της περιοχής
λειτουργίας τους. Δεύτερον, σε μετασχηματιστές σε σειρά, όπως ένας μετασχηματιστής 150/20 kV
που τροφοδοτείται από μετασχηματιστή 400/150 kV, το ΣΑΤΥΦ του ανώτερου επιπέδου τάσης πρέπει
να δρα πρώτο, με μικρότερη χρονική καθυστέρηση, ώστε να αποφεύγονται περιττοί χειρισμοί στο
κατώτερο επίπεδο. Επιπλέον, οι ροές άεργης ισχύος στις γραμμές διασύνδεσης με γειτονικά
συστήματα διατηρούνται συνήθως μικρές, ώστε κάθε ΔΣΜΗΕ να καλύπτει τις ανάγκες της περιοχής
ελέγχου του, ενώ η ανταλλαγή άεργης ισχύος στα σημεία σύνδεσης του ΣΜΗΕ με τα
δίκτυα διανομής αφορά και τους δύο διαχειριστές και εξετάζεται στο Κεφάλαιο 7.

\subsection{Ευστάθεια τάσης}
\label{sec:v-stability}

\subsubsection{Ορισμός και κατηγορίες}

Ευστάθεια τάσης (voltage stability) είναι η ικανότητα του συστήματος να διατηρεί αποδεκτές
τάσεις σε όλους τους ζυγούς μετά από μια διαταραχή, ξεκινώντας από δεδομένη κατάσταση
λειτουργίας. Η απώλειά της εκδηλώνεται συνήθως ως προοδευτική μείωση των τάσεων σε μια περιοχή
του συστήματος, η οποία μπορεί να καταλήξει σε \emph{κατάρρευση τάσης} (voltage collapse) και
σε διακοπή της τροφοδότησης μεγάλου μέρους της ζήτησης. Εμφανίζεται κυρίως σε φορτισμένα
συστήματα, όπου η ισχύς μεταφέρεται σε μεγάλες αποστάσεις από τους σταθμούς παραγωγής προς τα
κέντρα κατανάλωσης, και οφείλεται στο ότι το δίκτυο δεν μπορεί να καλύψει τη ζήτηση ισχύος, και
ιδίως άεργης ισχύος, των φορτίων. Διακρίνεται από την ευστάθεια γωνίας δρομέα (rotor angle
stability), η οποία αφορά τη
διατήρηση του συγχρονισμού των σύγχρονων μηχανών, αν και στην πράξη τα δύο φαινόμενα συχνά
συνυπάρχουν.

Η ευστάθεια τάσης κατατάσσεται με δύο κριτήρια. Ως προς το μέγεθος της διαταραχής, διακρίνεται
η ευστάθεια σε \emph{μεγάλες διαταραχές}, όπως βραχυκυκλώματα και απώλειες γραμμών ή μονάδων,
και η ευστάθεια σε \emph{μικρές διαταραχές}, όπως οι σταδιακές μεταβολές της ζήτησης. Ως προς
τη χρονική κλίμακα, η αστάθεια είναι \emph{βραχυπρόθεσμη}, με διάρκεια λίγων δευτερολέπτων,
όταν οφείλεται σε ταχέως αποκρινόμενα φορτία, όπως οι επαγωγικοί κινητήρες, ή σε μετατροπείς,
και \emph{μακροπρόθεσμη}, με διάρκεια λεπτών, όταν οφείλεται σε βραδύτερους μηχανισμούς, όπως
τα ΣΑΤΥΦ, τα θερμοστατικά ελεγχόμενα φορτία και οι περιοριστές υπερδιέγερσης.

\subsubsection{Μέγιστη μεταφερόμενη ισχύς και καμπύλες P--V}

Το φαινόμενο εξηγείται με το σύστημα του Σχήματος~\ref{fig:vq} για $R = 0$. Ο ζυγός 1
παριστάνει το υπόλοιπο σύστημα ως πηγή σταθερής τάσης $E$, και ο ζυγός 2, με τάση $V$,
τροφοδοτεί φορτίο $P + jQ$ με σταθερό συντελεστή ισχύος, δηλαδή $Q = P\tan\varphi$. Το
τετράγωνο του μέτρου της \eqref{eq:v1phasor} δίνει $E^2 = (V + XQ/V)^2 + (XP/V)^2$ ή,
ισοδύναμα,

\begin{equation}
V^4 + \left(2QX - E^2\right) V^2 + X^2\left(P^2 + Q^2\right) = 0
\label{eq:pv}
\end{equation}

Η \eqref{eq:pv} είναι δευτεροβάθμια ως προς $V^2$. Για κάθε τιμή της $P$ μέχρι ένα μέγιστο
έχει δύο λύσεις, δηλαδή δύο σημεία λειτουργίας: ένα με υψηλή τάση και μικρό ρεύμα, στο οποίο
λειτουργεί κανονικά το σύστημα, και ένα με χαμηλή τάση και μεγάλο ρεύμα. Καθώς η $P$ αυξάνεται,
οι δύο λύσεις πλησιάζουν και συμπίπτουν όταν μηδενιστεί η διακρίνουσα, στη μέγιστη ισχύ που
μπορεί να μεταφερθεί:

\begin{equation}
P_{\max} = \frac{E^2 \cos\varphi}{2X\left(1 + \sin\varphi\right)}
\label{eq:pmax}
\end{equation}

Για μοναδιαίο συντελεστή ισχύος $P_{\max} = E^2/(2X)$ και η αντίστοιχη τάση είναι
$E/\sqrt{2}$, δηλαδή περίπου 0,71 της τάσης της πηγής. Η μέγιστη ισχύς είναι αντιστρόφως
ανάλογη της αντίδρασης $X$, οπότε η απώλεια μιας γραμμής, η οποία αυξάνει την ισοδύναμη
αντίδραση, τη μειώνει.

Η γραφική παράσταση της $V$ ως συνάρτησης της $P$ ονομάζεται καμπύλη $P$--$V$ ή, λόγω του
σχήματός της, καμπύλη μύτης (nose curve). Το Σχήμα~\ref{fig:pv} δείχνει τις καμπύλες σε
κανονικοποιημένη μορφή για διάφορους συντελεστές ισχύος. Το επάνω τμήμα κάθε καμπύλης είναι η
περιοχή κανονικής λειτουργίας, ενώ στο κάτω τμήμα η λειτουργία είναι ασταθής για φορτία που
τείνουν να αποκαταστήσουν την ισχύ τους, όπως εξηγείται στη συνέχεια. Η απόσταση του σημείου
λειτουργίας από τη μύτη της καμπύλης, το \emph{περιθώριο φόρτισης} (loadability margin),
είναι το βασικό μέτρο της εγγύτητας στην αστάθεια τάσης.

\begin{figure}[!htbp]
\centering
\includegraphics[width=\linewidth]{figures/pv-curves.svg}
\caption{Καμπύλες $P$--$V$ του συστήματος δύο ζυγών για διάφορες τιμές του $\tan\varphi$ του
φορτίου, σε κανονικοποιημένη μορφή. Θετικές τιμές αντιστοιχούν σε επαγωγικό φορτίο και
αρνητικές σε φορτίο με αντιστάθμιση μεγαλύτερη από την άεργη ζήτησή του. Οι συνεχείς γραμμές
είναι τα επάνω τμήματα, στα οποία λειτουργεί κανονικά το σύστημα, και οι διακεκομμένες τα
κάτω. Οι κουκκίδες είναι τα κρίσιμα σημεία μέγιστης μεταφερόμενης ισχύος.}
\label{fig:pv}
\end{figure}

Από το σχήμα προκύπτουν δύο συμπεράσματα. Πρώτον, τα επαγωγικά φορτία μειώνουν τη μέγιστη
ισχύ, ενώ η αντιστάθμιση της άεργης ζήτησής τους την αυξάνει. Δεύτερον, με την αντιστάθμιση
η τάση στη μύτη της καμπύλης, η \emph{κρίσιμη τάση}, αυξάνεται και πλησιάζει τις κανονικές
τιμές λειτουργίας. Σε ισχυρά αντισταθμισμένο σύστημα η τάση μπορεί να παραμένει αποδεκτή μέχρι
πολύ κοντά στο όριο, οπότε το μέτρο της τάσης δεν αποτελεί από μόνο του αξιόπιστο δείκτη της
εγγύτητας στην κατάρρευση.

\subsubsection{Καμπύλες Q--V και περιθώριο άεργης ισχύος}

Συμπληρωματική εικόνα δίνει η καμπύλη $Q$--$V$ ενός ζυγού: η άεργη ισχύς $Q_c$ που πρέπει να
εγχέει ένας πλασματικός αντισταθμιστής στον ζυγό, ώστε η τάση του να είναι $V$, για δεδομένη
ζήτηση $P + jQ$. Για το σύστημα δύο ζυγών, με απαλοιφή της $\delta$ από την
\eqref{eq:pq-line},

\begin{equation}
Q_c = Q + \frac{V^2}{X} - \sqrt{\left(\frac{EV}{X}\right)^2 - P^2}
\label{eq:qv}
\end{equation}

Οι καμπύλες του Σχήματος~\ref{fig:qv} παρουσιάζουν ελάχιστο. Δεξιά του ελαχίστου η αύξηση της
εγχεόμενης άεργης ισχύος αυξάνει την τάση ($\mathrm{d}Q_c/\mathrm{d}V > 0$), όπως στην
κανονική λειτουργία. Αριστερά του ελαχίστου η σχέση αντιστρέφεται, οπότε η ρύθμιση της τάσης
με έγχυση άεργης ισχύος οδηγεί σε αστάθεια. Το σημείο λειτουργίας χωρίς αντιστάθμιση είναι η
τομή του δεξιού κλάδου με τον άξονα $Q_c = 0$, και η απόσταση του ελαχίστου από τον άξονα
αυτόν είναι το \emph{περιθώριο άεργης ισχύος} του ζυγού: η επιπλέον άεργη ζήτηση που μπορεί να
καλυφθεί πριν από την αστάθεια. Όταν το ελάχιστο βρίσκεται πάνω από τον άξονα, το σύστημα δεν
μπορεί να λειτουργήσει χωρίς αντιστάθμιση. Στην πράξη οι καμπύλες $Q$--$V$ υπολογίζονται με
διαδοχικές επιλύσεις ροής ισχύος, στις οποίες ο εξεταζόμενος ζυγός θεωρείται ζυγός ελεγχόμενης
τάσης χωρίς όρια άεργης ισχύος.

\begin{figure}[!htbp]
\centering
\includegraphics[width=\linewidth]{figures/qv-curves.svg}
\caption{Καμπύλες $Q$--$V$ στον ζυγό φορτίου του συστήματος δύο ζυγών, σε κανονικοποιημένη
μορφή, για φορτίο με $\tan\varphi$ = 0,2 και διάφορες τιμές της ενεργής ισχύος. Οι κουκκίδες
είναι τα ελάχιστα, που χωρίζουν τον ευσταθή δεξιό κλάδο από τον ασταθή αριστερό
(διακεκομμένη γραμμή), και οι κύκλοι τα σημεία λειτουργίας χωρίς αντιστάθμιση.}
\label{fig:qv}
\end{figure}

\subsubsection{Ο ρόλος των φορτίων}
\label{sec:load-role}

Η αστάθεια τάσης οφείλεται κυρίως στη δυναμική των φορτίων. Η ισχύς ενός φορτίου εξαρτάται
από την τάση, συνήθως σύμφωνα με το εκθετικό μοντέλο

\begin{equation}
P = P_0 \left(\frac{V}{V_0}\right)^{\alpha}, \qquad
Q = Q_0 \left(\frac{V}{V_0}\right)^{\beta}
\label{eq:loadexp}
\end{equation}

όπου $P_0$, $Q_0$ η ισχύς στην ονομαστική τάση $V_0$. Για $\alpha = 0$ το φορτίο είναι σταθερής
ισχύος, για $\alpha = 1$ σταθερού ρεύματος και για $\alpha = 2$ σταθερής σύνθετης αντίστασης·
για τα συνήθη φορτία ισχύει $0 < \alpha < 2$ και $\beta > \alpha$. Το σημείο λειτουργίας είναι
η τομή της καμπύλης $P$--$V$ του δικτύου με τη χαρακτηριστική του φορτίου.

Η εξάρτηση αυτή αφορά όμως μόνο τα πρώτα δευτερόλεπτα μετά από μια μεταβολή της τάσης.
Διάφοροι μηχανισμοί τείνουν στη συνέχεια να επαναφέρουν την ισχύ των φορτίων στην αρχική της
τιμή: τα ΣΑΤΥΦ των μετασχηματιστών διανομής επαναφέρουν σε μερικά λεπτά την τάση
στην πλευρά του φορτίου (Ενότητα «\nameref{sec:oltc}»), τα θερμοστατικά ελεγχόμενα φορτία
παραμένουν περισσότερο σε λειτουργία και οι επαγωγικοί κινητήρες αυξάνουν την ολίσθησή τους.
Μακροπρόθεσμα η ζήτηση συμπεριφέρεται επομένως ως ζήτηση σταθερής ισχύος.

Το Σχήμα~\ref{fig:load-char} δείχνει τη διαδικασία που οδηγεί σε κατάρρευση. Πριν από το
απρόβλεπτο συμβάν το σύστημα λειτουργεί στο σημείο A. Η απώλεια μιας γραμμής αυξάνει την
αντίδραση του δικτύου και συρρικνώνει την καμπύλη $P$--$V$. Αμέσως μετά το συμβάν, η τάση και
μαζί της η ζήτηση μειώνονται και το σύστημα μεταβαίνει στο σημείο B, στην τομή της νέας καμπύλης
με τη χαρακτηριστική του φορτίου. Στη συνέχεια το ΣΑΤΥΦ αυξάνει διαδοχικά τον λόγο
μετασχηματισμού, ώστε να επαναφέρει την τάση του φορτίου. Κάθε χειρισμός μετατοπίζει τη
χαρακτηριστική προς μεγαλύτερη ισχύ και το σημείο λειτουργίας προς τη μύτη της νέας καμπύλης.
Αν η αρχική ζήτηση $P_0$ είναι μικρότερη από τη νέα μέγιστη ισχύ, η αποκατάσταση ολοκληρώνεται
σε ένα νέο σημείο ισορροπίας με χαμηλότερη τάση στο δίκτυο μεταφοράς. Αν είναι μεγαλύτερη,
όπως στο σχήμα, δεν υπάρχει σημείο ισορροπίας με αποκατεστημένη ζήτηση. Μετά τη μύτη κάθε
χειρισμός του ΣΑΤΥΦ έχει το αντίθετο αποτέλεσμα: το σημείο λειτουργίας περνά στο κάτω τμήμα
της καμπύλης, οι τάσεις μειώνονται προοδευτικά, το ρεύμα αυξάνεται και η πορεία επιταχύνεται
όταν οι γεννήτριες φθάνουν στο όριο διέγερσης και ο OEL περιορίζει την άεργη ισχύ τους. Η
χρονική εξέλιξη του φαινομένου προσομοιώνεται στα Παραδείγματα.

\begin{figure}[!htbp]
\centering
\includegraphics[width=\linewidth]{figures/load-restoration.svg}
\caption{Αποκατάσταση του φορτίου από το ΣΑΤΥΦ μετά από απρόβλεπτο συμβάν, σε
ανά μονάδα τιμές. Οι καμπύλες $P$--$V$ αντιστοιχούν στο δίκτυο πριν και μετά την απώλεια μιας
γραμμής και οι λεπτές καμπύλες στη χαρακτηριστική του φορτίου για διαδοχικές λήψεις του
ΣΑΤΥΦ. Το φορτίο έχει $\alpha = \beta$ = 1 και, λόγω αντιστάθμισης, $\tan\varphi$ = −0,2. Οι
κουκκίδες είναι τα διαδοχικά σημεία ισορροπίας. Η αρχική ζήτηση $P_0$ υπερβαίνει τη μέγιστη ισχύ
του δικτύου μετά το συμβάν, οπότε η αποκατάστασή της οδηγεί το σημείο λειτουργίας στη μύτη της
καμπύλης.}
\label{fig:load-char}
\end{figure}

\subsubsection{Ανάλυση ευαισθησίας και δείκτες εγγύτητας}

Σε δίκτυο πολλών ζυγών οι καμπύλες $P$--$V$ και $Q$--$V$ υπολογίζονται με επαναλαμβανόμενες
ροές ισχύος, ενώ η εγγύτητα στο όριο εκτιμάται και από τη γραμμικοποίηση των εξισώσεων ροής
ισχύος γύρω από το σημείο λειτουργίας:

\begin{equation}
\begin{bmatrix} \Delta\mathbf{P} \\ \Delta\mathbf{Q} \end{bmatrix} =
\begin{bmatrix} \mathbf{J}_{P\theta} & \mathbf{J}_{PV} \\
                \mathbf{J}_{Q\theta} & \mathbf{J}_{QV} \end{bmatrix}
\begin{bmatrix} \Delta\boldsymbol{\theta} \\ \Delta\mathbf{V} \end{bmatrix}
\label{eq:pf-jacobian}
\end{equation}

όπου ο πίνακας είναι ο ιακωβιανός της ροής ισχύος, με παραγώγους της ίδιας μορφής με τις
παραγώγους των εγχύσεων στην εκτίμηση κατάστασης (Κεφάλαιο 1, Ενότητα «\nameref{sec:se}»).
Για σταθερή ενεργή ισχύ, $\Delta\mathbf{P} = \mathbf{0}$, η απαλοιφή των γωνιών δίνει

\begin{equation}
\Delta\mathbf{Q} = \mathbf{J}_R\,\Delta\mathbf{V}, \qquad
\mathbf{J}_R = \mathbf{J}_{QV} - \mathbf{J}_{Q\theta}\,\mathbf{J}_{P\theta}^{-1}\,\mathbf{J}_{PV}
\label{eq:jr}
\end{equation}

Ο μειωμένος ιακωβιανός $\mathbf{J}_R$ συνδέει τις μεταβολές της άεργης ισχύος με τις
μεταβολές των τάσεων και γενικεύει την κλίση των καμπυλών $Q$--$V$. Τα διαγώνια στοιχεία του
$\mathbf{J}_R^{-1}$ είναι οι ευαισθησίες τάσης--άεργης ισχύος των ζυγών: θετικές τιμές
αντιστοιχούν σε ευσταθή λειτουργία και μεγάλες τιμές σε ασθενείς ζυγούς. Καθώς η φόρτιση
αυξάνεται, η μικρότερη ιδιοτιμή του $\mathbf{J}_R$ μειώνεται και μηδενίζεται στη μύτη της
καμπύλης $P$--$V$, όπου ο ιακωβιανός γίνεται ιδιάζων. Η \emph{ανάλυση ιδιομορφών} (modal
analysis) εξετάζει τις μικρότερες ιδιοτιμές και τα αντίστοιχα ιδιοδιανύσματα: οι
\emph{συντελεστές συμμετοχής} των ζυγών δείχνουν ποιες περιοχές συμμετέχουν περισσότερο σε
κάθε ασθενή ιδιομορφή και επομένως πού είναι αποτελεσματικότερη η αντιστάθμιση.

Οι ΔΣΜΗΕ υπολογίζουν το περιθώριο φόρτισης με \emph{ροή ισχύος συνέχισης} (continuation
power flow), η οποία ακολουθεί την καμπύλη $P$--$V$ έως τη μύτη, για την κατάσταση N και για τα
απρόβλεπτα συμβάντα του καταλόγου, και απαιτούν ένα ελάχιστο περιθώριο να διατηρείται και μετά
από κάθε συμβάν. Η συνήθης ροή ισχύος αποκλίνει κοντά στη μύτη, ακριβώς επειδή ο ιακωβιανός
γίνεται ιδιάζων.

\subsubsection{Μέτρα πρόληψης και αντιμετώπισης}

Τα \emph{προληπτικά} μέτρα στοχεύουν στη διατήρηση επαρκούς περιθωρίου φόρτισης: εφεδρεία
άεργης ισχύος, ιδίως ταχείας, κοντά στα κέντρα κατανάλωσης, έγκαιρη ζεύξη πυκνωτών, κατάλληλες
τάσεις αναφοράς των μονάδων και, όπου εφαρμόζεται, δευτερεύουσα ρύθμιση τάσης. Τα
\emph{διορθωτικά} μέτρα εφαρμόζονται όταν το σύστημα πλησιάζει το όριο: αναστολή των χειρισμών
ή μείωση της τάσης αναφοράς των ΣΑΤΥΦ, ώστε να διακοπεί η αποκατάσταση του
φορτίου, ταχεία ζεύξη πυκνωτών, βραχυχρόνια υπερφόρτιση των γεννητριών και, ως έσχατο μέτρο του
σχεδίου άμυνας του συστήματος, αυτόματη αποσύνδεση ζήτησης λόγω χαμηλής τάσης (undervoltage
load shedding).

\section{Γρήγορη έγχυση άεργου ρεύματος}
\label{sec:frt}

Κατά τη διάρκεια ενός βραχυκυκλώματος η τάση μειώνεται απότομα σε μια περιοχή του συστήματος,
η έκταση της οποίας εξαρτάται από την ισχύ βραχυκύκλωσης του δικτύου. Η κύρια προστασία
απομονώνει το σφάλμα σε χρόνο της τάξης των 100 ms, αλλά οι επιπτώσεις του μπορεί να είναι
πολύ μεγαλύτερες. Για τις μονάδες ηλεκτροπαραγωγής ορίζονται δύο διακριτές απαιτήσεις, οι
οποίες συχνά συγχέονται (Πίνακας~\ref{tab:frt-vs-q}):

\begin{itemize}
\item Η \textbf{αδιάλειπτη λειτουργία υπό σφάλμα} (fault ride-through, FRT) είναι
      απαίτηση \emph{αντοχής}: η μονάδα δεν πρέπει να αποσυνδεθεί κατά τη διάρκεια της
      βύθισης και πρέπει να επανέλθει σε κανονική λειτουργία μετά την απομόνωση του σφάλματος.
      Αποτρέπει την απώλεια παραγωγής, αλλά δεν απαιτεί από μόνη της στήριξη της τάσης.
\item Η \textbf{γρήγορη έγχυση άεργου ρεύματος} (fast reactive current injection) είναι
      απαίτηση \emph{στήριξης}: η μονάδα πρέπει να εγχέει κατά τη διάρκεια της βύθισης άεργο
      ρεύμα που συγκρατεί την τάση. Αυτή είναι η επικουρική υπηρεσία του
      Πίνακα~\ref{tab:nfas} και προϋποθέτει την αδιάλειπτη λειτουργία.
\end{itemize}

\begin{table}[!htbp]
\centering
\begin{tabular}{lll}
\hline
 & Αδιάλειπτη λειτουργία & Έγχυση άεργου ρεύματος \\
\hline
είδος απαίτησης & αντοχή & στήριξη \\
τι ζητείται & να μην αποσυνδεθεί & άεργο ρεύμα στη βύθιση \\
μορφή & καμπύλη τάσης--χρόνου & $\Delta I_q = k\,\Delta U$ \\
σκοπός & να μη χαθεί παραγωγή & συγκράτηση της τάσης \\
σχέση με την υπηρεσία & προϋπόθεση & η ίδια η υπηρεσία \\
\hline
\end{tabular}
\caption{Αδιάλειπτη λειτουργία υπό σφάλμα και γρήγορη έγχυση άεργου ρεύματος.}
\label{tab:frt-vs-q}
\end{table}

Η διάκριση έχει σημασία κυρίως για τις μονάδες με μετατροπείς. Η σύγχρονη μηχανή καλύπτει και
τις δύο απαιτήσεις εγγενώς, ενώ ένας μετατροπέας μπορεί να παραμένει συνδεδεμένος χωρίς να
στηρίζει την τάση.

\subsection{Απόκριση των σύγχρονων μηχανών}
\label{sec:sm-response}

Αμέσως μετά από μια μεταβολή της τάσης στους ακροδέκτες της, η σύγχρονη μηχανή συμπεριφέρεται
ως πηγή τάσης πίσω από την υπομεταβατική αντίδραση $X_d^{\prime\prime}$, τυπικά 0,15--0,25 p.u.
Μια βύθιση της τάσης κατά $\Delta V$ προκαλεί επομένως ακαριαία μεταβολή του ρεύματος

\begin{equation}
\Delta I \approx \frac{\Delta V}{X_d^{\prime\prime}}
\label{eq:sm-dI}
\end{equation}

σε ανά μονάδα τιμές, χωρίς καμία ενέργεια ελέγχου. Σε τριφασικό βραχυκύκλωμα στους ακροδέκτες
το αρχικό ρεύμα φθάνει τις 4 έως 7 φορές το ονομαστικό και στη συνέχεια μειώνεται, καθώς η
μηχανή περνά από την υπομεταβατική στη μεταβατική και στη μόνιμη κατάσταση. Στο μεταξύ ο
ρυθμιστής τάσης αυξάνει τη διέγερση στη μέγιστη τιμή της (field forcing) και συγκρατεί την
τάση. Επειδή η σύνθετη αντίσταση του δικτύου και του σφάλματος είναι κυρίως επαγωγική, το ρεύμα
αυτό είναι κατά κύριο λόγο άεργο. Η σύγχρονη μηχανή παραμένει σε λειτουργία, εφόσον διατηρεί
τον συγχρονισμό της, και ταυτόχρονα στηρίζει την τάση.

\subsection{Αδιάλειπτη λειτουργία υπό σφάλμα}
\label{sec:frt-req}

Οι μονάδες που συνδέονται μέσω μετατροπέων συμπεριφέρονται διαφορετικά. Το ρεύμα του
μετατροπέα μπορεί να υπερβεί το ονομαστικό μόνο κατά 10--20\%, γιατί οι ημιαγωγικοί διακόπτες
έχουν πολύ μικρή θερμική χωρητικότητα και δεν αντέχουν υπερφόρτιση ούτε για λίγα ms. Επιπλέον,
το ρεύμα καθορίζεται από τον έλεγχο: ένας μετατροπέας που ακολουθεί το δίκτυο (grid-following)
μετρά τη γωνία της τάσης με βρόχο κλειδωμένης φάσης (Phase-Locked Loop, PLL) και εγχέει το
ρεύμα που ορίζουν οι ελεγκτές του. Οι πρώτες ανεμογεννήτριες αποσυνδέονταν σε κάθε βύθιση της
τάσης, για να προστατευθούν. Με την αύξηση της αιολικής παραγωγής, η μαζική αποσύνδεση μετά από
ένα σφάλμα στο ΣΜΗΕ έγινε απειλή για την ευστάθεια του συστήματος: ένα τοπικό
σφάλμα μπορούσε να προκαλέσει απώλεια παραγωγής μεγαλύτερη από το συμβάν αναφοράς, για το οποίο
διαστασιολογείται η εφεδρεία διατήρησης συχνότητας (Κεφάλαιο 1, Ενότητα «\nameref{sec:n1}»).
Οι κώδικες σύνδεσης εισήγαγαν γι' αυτό την απαίτηση αδιάλειπτης λειτουργίας υπό σφάλμα.

Η απαίτηση διατυπώνεται ως καμπύλη τάσης--χρόνου στο σημείο σύνδεσης (Σχήμα~\ref{fig:frt}).
Η μονάδα πρέπει να παραμένει συνδεδεμένη και σε ευσταθή λειτουργία όσο η τάση βρίσκεται πάνω
από την καμπύλη, ενώ επιτρέπεται να αποσυνδεθεί όταν η τάση πέσει κάτω από αυτήν. Ο RfG
κατατάσσει τις μονάδες σε τέσσερις τύπους, Α έως Δ, ανάλογα με την ισχύ τους και την τάση
σύνδεσης· στον τύπο Δ ανήκουν οι μονάδες που συνδέονται σε τάση 110 kV και άνω. Για κάθε τύπο
ορίζει τη μορφή της καμπύλης και τα εύρη των παραμέτρων της, ενώ τις τιμές καθορίζει ο ΔΣΜΗΕ. Για
μονάδες τύπου Δ που συνδέονται μέσω μετατροπέων, η καμπύλη μπορεί να απαιτεί αδιάλειπτη
λειτουργία ακόμη και με μηδενική τάση για 140--150 ms, χρόνο που καλύπτει την απομόνωση του
σφάλματος από την κύρια προστασία, και στη συνέχεια όσο η τάση ανακάμπτει πάνω από μια ευθεία
που φθάνει τα 0,85 p.u. σε 1,5--3 s. Πολλοί κώδικες ορίζουν αντίστοιχες απαιτήσεις και για τις
υπερτάσεις που ακολουθούν την απομόνωση του σφάλματος.

\begin{figure}[!htbp]
\centering
\includegraphics[width=\linewidth]{figures/frt-profile.svg}
\caption{Ενδεικτική καμπύλη αδιάλειπτης λειτουργίας υπό σφάλμα για μονάδα που συνδέεται μέσω
μετατροπέα και μια τυπική τροχιά της τάσης στο σημείο σύνδεσης, η οποία πέφτει στα 0,25 p.u. κατά το
σφάλμα και ανακάμπτει μετά την απομόνωσή του στα 120 ms. Η μονάδα πρέπει να παραμένει σε
λειτουργία όσο η τάση βρίσκεται στη σκιασμένη περιοχή πάνω από την καμπύλη.}
\label{fig:frt}
\end{figure}

Η απαίτηση αφορά μόνο το αν η μονάδα παραμένει συνδεδεμένη και επανέρχεται σε κανονική
λειτουργία. Ένας μετατροπέας μπορεί να την ικανοποιεί ακόμη και με μηδενικό ρεύμα κατά τη
διάρκεια της βύθισης, αν αναστείλει προσωρινά τη λειτουργία των διακοπτών του (blocking)· στην
περίπτωση αυτή δεν στηρίζει την τάση. Για να παραμείνει σε λειτουργία, η μονάδα πρέπει να
διαχειριστεί την ενέργεια της πρωτογενούς πηγής που δεν μπορεί να εγχυθεί στο δίκτυο όσο η
τάση είναι χαμηλή: η ενέργεια αυτή αποθηκεύεται προσωρινά ως κινητική ενέργεια στον δρομέα της
ανεμογεννήτριας, καταναλώνεται σε αντιστάτη πέδησης στον ζυγό ΣΡ ή, στα φωτοβολταϊκά,
περιορίζεται με μετακίνηση του σημείου λειτουργίας των πλαισίων. Μετά την απομόνωση του
σφάλματος η μονάδα πρέπει να επαναφέρει την ενεργή ισχύ της σε χρόνο που ορίζει ο ΔΣΜΗΕ. Η
γρήγορη επαναφορά περιορίζει τη διαταραχή της συχνότητας, αλλά σε ασθενή δίκτυα επιβαρύνει την
ανάκτηση της τάσης.

\subsection{Στήριξη της τάσης με έγχυση άεργου ρεύματος}
\label{sec:kfactor}

Η επικουρική υπηρεσία είναι η ενεργή στήριξη της τάσης κατά τη διάρκεια της βύθισης. Οι
κώδικες σύνδεσης απαιτούν από τις μονάδες με μετατροπείς να εγχέουν άεργο ρεύμα ανάλογο της
απόκλισης της τάσης, το οποίο ο RfG ονομάζει \emph{ταχύ ρεύμα σφάλματος} (fast fault current):

\begin{equation}
\Delta I_q = k\,\Delta U
\label{eq:kfactor}
\end{equation}

όπου $\Delta U = U_0 - U$ η βύθιση της τάσης ως προς την τιμή $U_0$ πριν από το σφάλμα, σε ανά
μονάδα της ονομαστικής τάσης, $\Delta I_q$ η αύξηση του χωρητικού άεργου ρεύματος, σε ανά
μονάδα του ονομαστικού ρεύματος, και $k$ ο συντελεστής ενίσχυσης, με τυπική τιμή 2: κάθε 1\%
βύθισης της τάσης απαιτεί 2\% επιπλέον άεργο ρεύμα. Σε ορισμένους κώδικες η χαρακτηριστική
έχει νεκρή ζώνη γύρω από την ονομαστική τάση, ώστε η στήριξη να ενεργοποιείται μόνο σε
πραγματικές βυθίσεις. Η απόκριση πρέπει να ολοκληρώνεται σε μερικές δεκάδες ms. Σε υπερτάσεις,
όπου $\Delta U < 0$, η ίδια χαρακτηριστική δίνει επαγωγικό ρεύμα. Η σύγκριση με την
\eqref{eq:sm-dI} δείχνει ότι η σύγχρονη μηχανή αποκρίνεται με ισοδύναμο συντελεστή
$1/X_d^{\prime\prime}$, δηλαδή 4 έως 7, και χωρίς το όριο ρεύματος του μετατροπέα
(Σχήμα~\ref{fig:iq}).

\begin{figure}[!htbp]
\centering
\includegraphics[width=\linewidth]{figures/frt-iq.svg}
\caption{Χαρακτηριστική έγχυσης άεργου ρεύματος μετατροπέα με $k$ = 2, η οποία περιορίζεται
από το μέγιστο ρεύμα του, και εγγενής απόκριση σύγχρονης μηχανής με $X_d^{\prime\prime}$ = 0,2
p.u.}
\label{fig:iq}
\end{figure}

Επειδή το συνολικό ρεύμα του μετατροπέα είναι περιορισμένο,

\begin{equation}
I_p^2 + I_q^2 \le I_{\max}^2
\label{eq:ilimit}
\end{equation}

όπου $I_p$ και $I_q$ η ενεργή και η άεργη συνιστώσα του ρεύματος, κατά τη διάρκεια της βύθισης
δίνεται προτεραιότητα στο άεργο ρεύμα και το ενεργό ρεύμα μειώνεται όσο χρειάζεται. Η ενεργή
ισχύς μειώνεται επομένως για δύο λόγους, επειδή μειώνονται τόσο η τάση όσο και το ενεργό ρεύμα.
Η άεργη ισχύς που εγχέεται είναι $U I_q$, δηλαδή μικρότερη από εκείνη που θα έδινε το ίδιο
ρεύμα σε ονομαστική τάση. Η στήριξη είναι αποτελεσματικότερη σε ασθενή δίκτυα, γιατί, σύμφωνα με
την \eqref{eq:dv-ssc}, η ίδια έγχυση μεταβάλλει την τάση τόσο περισσότερο όσο μικρότερη είναι η
ισχύς βραχυκύκλωσης.

Τα περισσότερα σφάλματα είναι μονοφασικά και προκαλούν ασύμμετρες βυθίσεις. Ένας μετατροπέας
που εγχέει μόνο ρεύμα θετικής ακολουθίας δεν περιορίζει την τάση αρνητικής ακολουθίας και
μπορεί να δυσχεράνει τη λειτουργία της προστασίας, η οποία χρησιμοποιεί συχνά μεγέθη αρνητικής
ακολουθίας. Γι' αυτό νεότεροι κώδικες απαιτούν και έγχυση ρεύματος αρνητικής ακολουθίας, κατ'
αναλογία με την \eqref{eq:kfactor}.

Ως κύρια λειτουργία τους, τη γρήγορη έγχυση άεργου ρεύματος παρέχουν οι STATCOM και οι
σύγχρονοι αντισταθμιστές. Ο STATCOM διατηρεί το ονομαστικό του ρεύμα σε πολύ χαμηλές τάσεις
(Ενότητα «\nameref{sec:facts}»), ενώ ο σύγχρονος αντισταθμιστής δίνει στους πρώτους κύκλους,
σύμφωνα με την \eqref{eq:sm-dI}, ρεύμα πολλαπλάσιο του ονομαστικού και συμβάλλει επιπλέον στο
ρεύμα βραχυκύκλωσης και στην αδράνεια. Η ταχεία στήριξη βοηθά και την ανάκτηση της τάσης μετά
το σφάλμα, όταν οι επαγωγικοί κινητήρες που επιβραδύνθηκαν απορροφούν μεγάλο άεργο ρεύμα για να
επιταχυνθούν, φαινόμενο που συνδέεται με τη βραχυπρόθεσμη αστάθεια τάσης (Ενότητα
«\nameref{sec:v-stability}»).

\section{Επανεκκίνηση και νησιδοποιημένη λειτουργία}
\label{sec:restoration}

Παρά τα μέτρα του σχεδίου άμυνας, ένα σύστημα μπορεί να μεταβεί σε κατάσταση γενικής διακοπής
(Κεφάλαιο 1, Ενότητα «\nameref{sec:states}»). Ο ΔΣΜΗΕ αποκαθιστά τότε την τροφοδότηση σύμφωνα με
το σχέδιο αποκατάστασης που προβλέπει ο Κανονισμός (ΕΕ) 2017/2196 για την κατάσταση έκτακτης
ανάγκης και την αποκατάσταση (Network Code on Emergency and Restoration, NC ER). Στόχος είναι η
ταχεία επανατροφοδότηση της ζήτησης χωρίς νέα κατάρρευση, η οποία θα επανέφερε τη διαδικασία
στην αρχή. Οι δύο υπηρεσίες της ενότητας, η ικανότητα επανεκκίνησης και η ικανότητα
νησιδοποιημένης λειτουργίας, είναι τα βασικά εργαλεία της αποκατάστασης.

\subsection{Στρατηγικές και στάδια της αποκατάστασης}
\label{sec:restoration-stages}

Η επανατροφοδότηση ξεκινά από μια πηγή τάσης, η οποία μπορεί να προέρχεται από δύο κατευθύνσεις:

\begin{itemize}
\item Στη στρατηγική \textbf{από πάνω προς τα κάτω} (top-down), η τάση προέρχεται από γειτονικά
      συστήματα μέσω των διασυνδέσεων. Η στρατηγική είναι ταχύτερη, αλλά προϋποθέτει ότι τα
      γειτονικά συστήματα λειτουργούν κανονικά και διαθέτουν επαρκή εφεδρεία.
\item Στη στρατηγική \textbf{από κάτω προς τα πάνω} (bottom-up), η τάση προέρχεται από μονάδες
      του ίδιου του συστήματος με ικανότητα επανεκκίνησης, οι οποίες σχηματίζουν νησίδες που
      στη συνέχεια συγχρονίζονται.
\end{itemize}

Στην πράξη οι δύο στρατηγικές συνδυάζονται. Η αποκατάσταση εξελίσσεται σε τρία στάδια
(Σχήμα~\ref{fig:restoration}):

\begin{enumerate}
\item \textbf{Προετοιμασία.} Ο ΔΣΜΗΕ προσδιορίζει την κατάσταση μετά τη διακοπή, δηλαδή ποια
      στοιχεία είναι διαθέσιμα και ποιες μονάδες μπορούν να εκκινήσουν, ανοίγει τους
      διακόπτες, ώστε το δίκτυο να ηλεκτρίζεται τμηματικά, και επιλέγει τις διαδρομές
      ηλέκτρισης.
\item \textbf{Αποκατάσταση του κορμού του δικτύου.} Οι μονάδες επανεκκίνησης ή οι διασυνδέσεις
      ηλεκτρίζουν επιλεγμένες διαδρομές υπερυψηλής τάσης προς μεγάλους σταθμούς, ώστε να
      τροφοδοτηθούν τα βοηθητικά τους και να επανεκκινήσουν. Σχηματίζονται νησίδες, στις οποίες
      η παραγωγή και το φορτίο αυξάνονται σταδιακά.
\item \textbf{Αποκατάσταση του φορτίου.} Οι νησίδες συγχρονίζονται μεταξύ τους και με τα
      γειτονικά συστήματα και το υπόλοιπο φορτίο επανατροφοδοτείται με τον ρυθμό που
      επιτρέπουν οι διαθέσιμες μονάδες.
\end{enumerate}

\begin{figure}[!htbp]
\centering
\includegraphics[width=\linewidth]{figures/restoration.svg}
\caption{Σχηματική απεικόνιση της αποκατάστασης από κάτω προς τα πάνω. Η μονάδα επανεκκίνησης
ηλεκτρίζει μια διαδρομή προς μεγάλο σταθμό (1), ο οποίος επανεκκινεί και σχηματίζει νησίδα με
φορτία που συνδέονται σε βήματα (2), και η νησίδα συγχρονίζεται τελικά με το γειτονικό σύστημα
(3).}
\label{fig:restoration}
\end{figure}

\subsection{Μονάδες επανεκκίνησης}
\label{sec:blackstart}

Ικανότητα επανεκκίνησης (black start capability) έχει μια μονάδα που μπορεί να εκκινήσει χωρίς
εξωτερική τροφοδότηση και να ηλεκτρίσει τμήμα του δικτύου. Η μονάδα πρέπει:

\begin{itemize}
\item να εκκινεί με δική της πηγή, όπως ντιζελογεννήτρια ή συσσωρευτές,
\item να ηλεκτρίζει γραμμές και μετασχηματιστές χωρίς φορτίο, άρα να απορροφά την άεργη ισχύ που
      παράγουν, εντός του ορίου υποδιέγερσής της (Ενότητα «\nameref{sec:sg}»),
\item να ρυθμίζει την τάση και τη συχνότητα με μικρό και μεταβαλλόμενο φορτίο και να αντέχει
      τις βηματικές μεταβολές του, και
\item να μπορεί να επαναλάβει την εκκίνηση, αν η πρώτη προσπάθεια αποτύχει.
\end{itemize}

Παραδοσιακά την ικανότητα αυτή παρέχουν υδροηλεκτρικοί σταθμοί, που χρειάζονται ελάχιστη ισχύ
για τα βοηθητικά τους, και αεριοστρόβιλοι με ντιζελογεννήτρια εκκίνησης· εξετάζονται επίσης
σταθμοί αποθήκευσης με μετατροπείς σχηματισμού δικτύου (grid-forming), οι οποίοι, αντίθετα
από τους μετατροπείς που ακολουθούν το δίκτυο, συμπεριφέρονται ως πηγή τάσης, όπως η σύγχρονη
μηχανή.
Ο RfG προβλέπει την ικανότητα επανεκκίνησης για μονάδες τύπου Γ και Δ χωρίς να την καθιστά
υποχρεωτική: παρέχεται κατόπιν αιτήματος του ΔΣΜΗΕ. Οι μονάδες επανεκκίνησης δοκιμάζονται
περιοδικά, ώστε να είναι βέβαιο ότι η ικανότητα θα είναι διαθέσιμη όταν απαιτηθεί.

\subsection{Τεχνικά προβλήματα της αποκατάστασης}
\label{sec:restoration-issues}

Κατά την αποκατάσταση το σύστημα είναι μικρό, ασθενές και ελαφρά φορτισμένο, οπότε τα
φαινόμενα των προηγούμενων ενοτήτων εμφανίζονται έντονα.

\textbf{Υπερτάσεις κατά την ηλέκτριση.} Μια γραμμή που ηλεκτρίζεται χωρίς φορτίο παράγει,
σύμφωνα με την \eqref{eq:qline} για $I \approx 0$, άεργη ισχύ $\omega C \ell V^2$, την οποία
πρέπει να απορροφήσουν οι λίγες μονάδες σε λειτουργία. Η τάση στο ανοικτό άκρο αυξάνεται
επιπλέον λόγω του φαινομένου Ferranti (Ενότητα «\nameref{sec:v-q}»). Για τον περιορισμό
των υπερτάσεων ηλεκτρίζονται σύντομες διαδρομές, συνδέονται πηνία, η τάση των μονάδων
ρυθμίζεται σε χαμηλές τιμές και συνδέεται νωρίς φορτίο. Η ζεύξη μετασχηματιστών χωρίς φορτίο
προκαλεί επιπλέον υψηλά ρεύματα μαγνήτισης (inrush currents) με έντονο αρμονικό περιεχόμενο, τα οποία, λόγω της
μικρής ισχύος βραχυκύκλωσης, μπορεί να διεγείρουν συντονισμούς και να προκαλέσουν προσωρινές
υπερτάσεις.

\textbf{Ρύθμιση της συχνότητας κατά τη σύνδεση φορτίων.} Η νησίδα έχει λίγες μονάδες και μικρή
αδράνεια. Σύμφωνα με τη σχέση του αρχικού ρυθμού μεταβολής της συχνότητας
(\hyperref[kef-03-theoria]{Κεφάλαιο 3}), η μεταβολή της συχνότητας για δεδομένη βηματική
μεταβολή της ζήτησης είναι αντιστρόφως ανάλογη της ισχύος των μονάδων σε λειτουργία, οπότε το
φορτίο συνδέεται σε μικρά βήματα. Τα φορτία που επανατροφοδοτούνται μετά από μακρά διακοπή
απορροφούν αρχικά περισσότερη ισχύ από την κανονική (cold load pickup), γιατί τα θερμοστατικά
ελεγχόμενα φορτία λειτουργούν τότε όλα ταυτόχρονα.

\textbf{Ασθενές δίκτυο και προστασία.} Η ισχύς βραχυκύκλωσης είναι μικρή, οπότε η προστασία
πρέπει να λειτουργεί με μικρά ρεύματα σφάλματος. Οι μονάδες με
μετατροπείς που ακολουθούν το δίκτυο συνδέονται μόνο όταν η νησίδα αποκτήσει επαρκή ισχύ
βραχυκύκλωσης.

\textbf{Συγχρονισμός και κλείσιμο βρόχων.} Για τη σύνδεση δύο νησίδων ή το κλείσιμο ενός
βρόχου, οι τάσεις στις δύο πλευρές του διακόπτη πρέπει να έχουν παραπλήσιο μέτρο, συχνότητα και
γωνία· τις συνθήκες αυτές ελέγχουν ηλεκτρονόμοι ελέγχου συγχρονισμού (synchro-check) πριν από
το κλείσιμο. Σε ανοικτό βρόχο μπορεί να εμφανιστεί μεγάλη διαφορά γωνίας (standing phase
angle), λόγω της ροής ενεργής ισχύος στην υπόλοιπη διαδρομή. Σύμφωνα με την
\eqref{eq:pq-line}, το κλείσιμο του διακόπτη προκαλεί τότε απότομη ροή ισχύος και καταπόνηση
των μονάδων, οπότε η διαφορά γωνίας περιορίζεται πρώτα με ανακατανομή της παραγωγής.

\subsection{Νησιδοποιημένη λειτουργία}
\label{sec:island}

Η ικανότητα νησιδοποιημένης λειτουργίας (island operation capability) απαιτεί από μια μονάδα
να ρυθμίζει τη συχνότητα και την τάση της νησίδας και να παραμένει σε ευσταθή λειτουργία μετά από
τις μεταβολές του φορτίου που συνοδεύουν τον σχηματισμό της. Ο ρυθμιστής στροφών πρέπει να
ανιχνεύει τη μετάβαση σε νησιδοποιημένη λειτουργία και να αλλάζει τρόπο ρύθμισης, για παράδειγμα σε
ισόχρονη ρύθμιση χωρίς στατισμό, αφού στη νησίδα δεν υπάρχει άλλη μονάδα να επαναφέρει τη
συχνότητα. Νησίδες σχηματίζονται και ακούσια, όταν το σύστημα διαχωρίζεται μετά από διαδοχικές
αποζεύξεις ή με ενέργεια του σχεδίου άμυνας. Οι μονάδες με ικανότητα νησιδοποιημένης λειτουργίας
συγκρατούν τότε τη συχνότητα και την τάση του τμήματος που απομονώθηκε και επιτρέπουν τη
γρήγορη επανασύνδεσή του.

Συγγενής ικανότητα είναι ο ταχύς επανασυγχρονισμός (quick re-synchronisation): μια μονάδα που
αποσυνδέεται από το δίκτυο παραμένει σε λειτουργία τροφοδοτώντας μόνο τα βοηθητικά της
(house load operation), ώστε να επανασυγχρονιστεί αμέσως μόλις επανέλθει η τάση, χωρίς τη
μακρά διαδικασία εκκίνησης ενός θερμικού σταθμού. Ο RfG προβλέπει την ικανότητα
νησιδοποιημένης λειτουργίας και τον ταχύ επανασυγχρονισμό για μονάδες τύπου Γ και Δ.

% ---------------------------------------------------------------------------
% BACKUP: εκτός δομής, για μελλοντική χρήση.
% Μελέτη περίπτωσης: γενική διακοπή στην Ιβηρική Χερσόνησο, 28 Απριλίου 2025.
% Η τελική έκθεση της ομάδας εμπειρογνωμόνων του ENTSO-E (Μάρτιος 2026) αποδίδει
% το συμβάν κυρίως σε υπερτάσεις και ανεπαρκή ρύθμιση τάσης, σε συνδυασμό με
% ταλαντώσεις που προηγήθηκαν. Πιθανή θέση: πλαίσιο μετά την ενότητα ρύθμισης
% τάσης ή στην ενότητα επανεκκίνησης.
% ---------------------------------------------------------------------------
